\documentclass[iicol,pdflatex,sn-basic,Numbered]{sn-jnl}

\usepackage{graphicx}
\usepackage{amsmath,amssymb,amsfonts}
\usepackage{amsthm}
\usepackage[title]{appendix}
\usepackage{xcolor}
\usepackage{booktabs}
\usepackage{multirow}

\theoremstyle{thmstyleone}
\newtheorem{theorem}{Theorem}
\newtheorem{proposition}[theorem]{Proposition}
\newtheorem{lemma}[theorem]{Lemma}
\newtheorem{corollary}[theorem]{Corollary}

\theoremstyle{thmstylethree}
\newtheorem{definition}[theorem]{Definition}
\newtheorem{assumption}[theorem]{Assumption}
\newtheorem{remark}[theorem]{Remark}

\raggedbottom

\begin{document}

\title[The Camouflage Ratio]{The Camouflage Ratio: Strategic Tie Formation and the\\
       Limits of Structural Detection}

\author*[1]{\fnm{Quang-Vinh} \sur{Dang}}\email{vinh.dq4@buv.edu.vn}

\author[1]{\fnm{Phuong-Lan} \sur{Nguyen}}

\affil*[1]{\orgname{British University Vietnam}, \orgaddress{\city{Hung Yen}, \country{Vietnam}}}

\abstract{Inferring an actor's attributes from those of its network neighbours rests on
homophily: neighbours are assumed informative. We ask what happens when an actor
\emph{chooses} its ties in order to defeat that inference, as automated accounts on
social platforms do when they form ties to genuine users so that their neighbourhood
resembles a genuine one. We formalise this with one structural parameter, the
\emph{camouflage ratio} $\rho$. Three results follow. First, camouflage is not
private: because a tie has two ends, edge-endpoint conservation forces the induced
contamination of the majority to be $\eta = c\rho$, where $c$ is the
minority-to-majority ratio, so prevalence acts as protection---the more common such
actors are, the less camouflage each needs. Second, a threshold
$\rho^\star = (1-D^{-1/2})/(1+c)$ in the mean neighbourhood size $D$ marks where
averaging over alters ceases to improve on attributes alone; we are explicit that a
crossover of this kind is not new, and that our contribution is the class-imbalanced
case. Third, selectively discarding suspicious ties does not straightforwardly repair
the inference, because it trades bias against the number of ties left to average
over. Measuring $\rho$ on $10{,}199$ expert-annotated Twitter accounts, every
recorded relation is one the network harms; yet the same accounts avoid one another
where ties are visible and converge where content is shared, so multiplexity rather
than aggregate density decides whether relational inference is usable.}

\keywords{Relational inference, Homophily, Strategic tie formation, Multiplexity, Social bots, Stochastic block models}

\maketitle

\section{Introduction}
\label{sec:intro}

One of the oldest uses of relational data is to infer something about an actor from
the actors it is connected to. The inference rests on \emph{homophily}: if ties are
formed preferentially between similar actors, then a neighbourhood is informative
about its centre, and averaging over it estimates the centre's attributes more
precisely than the centre's own noisy attributes do. Variants of this logic underpin
network autocorrelation models, peer-influence estimation, and relational
classification.

This paper asks what happens when an actor \emph{chooses} its ties in order to defeat
that inference. The setting where this is most consequential, and where labelled data
exists, is the detection of automated accounts on social platforms. Early generations
of such accounts formed dense, mutually connected clusters and were consequently easy
to identify collectively. Contemporary ones do the opposite: they form ties to large
numbers of genuine users so that their local neighbourhood is statistically
indistinguishable from a genuine actor's---a strategy known in that literature as
\emph{relation camouflage}~\citep{dou2020caregnn,cresci2020decade}. The strategy is
purposive and cheap, and it targets the inferential mechanism rather than any
particular detector. Detection systems have responded by discarding suspicious ties,
by treating similar and dissimilar neighbours through separate channels, or by holding
an actor's own attributes apart from its neighbours'~\citep{dou2020caregnn,he2025bothp},
and recent detectors combine all three~\citep{dang2026xhbot}. Those designs are
justified empirically, by ablation. What has been missing is an account of the
structural conditions under which the underlying inference works at all, which is what
we supply here.

We formalise it with one structural parameter, the \emph{camouflage ratio} $\rho$: the
share of a strategic actor's ties directed at the other class. Two questions then
have precise answers.

\begin{enumerate}
    \item \textbf{When is relational information worth using at all?} If camouflage
    is severe enough, averaging over neighbours must eventually be worse than
    ignoring the network and using attributes alone. Practitioners have no criterion
    for where that point lies, and---as we show---it depends on quantities they can
    already measure.
    \item \textbf{Does selectively discarding suspicious ties restore the
    inference?} Discarding ties is the standard remedy, and it is justified almost
    entirely by ablation: removing the step costs accuracy on a benchmark. That is
    not a statement about the conditions under which it works.
\end{enumerate}

Our analysis is conducted in a contextual stochastic block
model~\citep{deshpande2018contextual}---a stochastic block model with attributes on
the nodes---augmented with $\rho$ as an explicit strategic parameter, and we study
what neighbourhood averaging does to the separation between the two classes as $\rho$
grows.

\paragraph{What is already known.} We are explicit about priority from the outset.
For a \emph{balanced} symmetric CSBM, Ma et al.~\citep[Thm.~2]{ma2022homophily}
already established a closed-form threshold at which aggregation ceases to beat the
graph-free classifier, and our critical ratio reduces to their condition exactly when
the two classes are equal in size (Remark~\ref{rem:ma}); related crossover analyses
appear in~\citep{luan2023when,mao2023demystifying,fountoulakis2023graph}. The
existence of a crossover is therefore not our discovery, and we do not claim it.
Our contribution is the \emph{class-imbalanced adversarial} case, where the two
classes deflate by different amounts---which moves the decision boundary and makes
bot prevalence itself a robustness parameter---together with a closed form for how
far edge pruning shifts the threshold.

\paragraph{Contributions.}

\begin{itemize}
    \item \textbf{The edge-balance identity, and bot prevalence as a robustness
    parameter.} Camouflage edges have two ends, so a bot that points an edge at a
    human necessarily pollutes that human's neighbourhood too. We show this forces
    $\eta = c\rho$, where $c$ is the bot-to-human ratio, so that the total deflation
    is $1-(1+c)\rho$ (Proposition~\ref{prop:deflation}). The consequence is that
    bot prevalence enters the robustness threshold directly: as bots become more
    common, each one needs \emph{less} camouflage to break the detector. This is not
    a class-imbalance problem that resampling addresses.

    \item \textbf{A Bayes-error crossover valid under class imbalance.} When
    $c \neq 1$ the two classes deflate asymmetrically, the post-aggregation class
    midpoint moves, and a comparison against a fixed decision boundary is no longer
    valid---the objection Luan et al.~\citep{luan2023when} raise
    against~\citep{ma2022homophily}. We compare Bayes errors instead
    (Theorem~\ref{thm:error}), obtaining
    $\rho^\star = (1 - D^{-1/2})/(1+c)$ (Corollary~\ref{cor:critical}), which
    recovers~\citep{ma2022homophily} at $c=1$ and~\citep{baranwal2021graph} at
    $\rho=0$. Our error term also retains a mixture-variance contribution that a
    naive $\sigma^2/D$ calculation misses, and which shifts the threshold by $11.5\%$.

    \item \textbf{What discarding ties can and cannot buy.} Selectively removing
    suspicious ties is the standard remedy, and we show it is a trade rather than a
    free improvement: down-weighting ties non-uniformly also reduces the number of
    ties left to average over, so the benefit is bounded
    (Proposition~\ref{prop:deff}). We give the tolerance as a closed function of the
    discriminating power of whatever suspicion score is used, and bound that power
    from the distribution of attribute discrepancies, which removes the free parameter
    (Section~\ref{sec:achievable}). On the review networks we measure, a score based on
    attribute discrepancy alone moves the threshold by less than $0.01$ and does not
    restore the inference, while a flawless one would---so the operative quantity is
    the gap between achievable and ideal tie scoring.

    \item \textbf{Measurement on real networks.} We give an estimator for $\rho$,
    validate it against planted ground truth, and report it for several labelled
    populations. Camouflage differs by a factor approaching two across relation types
    recorded on the \emph{same} actors, which makes these per-relation quantities and
    implies that pooling relations averages a usable signal with an unusable one.
\end{itemize}

Our aim is not to propose a detector. It is to characterise a structural mechanism
that limits what ties can reveal, to show that its governing quantities are
measurable, and to report what they are on real networks. The practical consequence is
that ``should I use the network here?'' becomes a question that can be answered from
data before any model is fitted, and answered separately for each type of tie.



\section{Related Work}
\label{sec:related}

\subsection{Inferring attributes from alters}
\label{sec:related-alters}

The operation we analyse has a long history in network analysis under several names.
In the network autocorrelation model an actor's outcome is regressed on a weighted
average of its alters' outcomes, and Leenders~\citep{leenders2002modeling} makes the
point on which this paper turns: the weight matrix is a substantive modelling
commitment, not a technical detail, and the conclusions depend on it. Row-normalised
weights---the natural choice, and the one implicit in most applications---are exactly
the mean aggregation of Section~\ref{sec:model}. Everett and
Borgatti~\citep{everett2026alter} give the general matrix formulation of \emph{alter
composition}, the profile of categories an actor is exposed to through its alters,
and extend it to overlapping memberships; our camouflage ratio is a two-category
alter-composition measure, and the question we ask is what happens to it when one
category chooses its ties.

Two in-journal results bear directly on the reliability of this inference.
Roeling and Nicholls~\citep{roeling2020imputation} impute unobserved attributes from
network structure and find the procedure sensitive to model misspecification and to
missingness. V\"or\"os et al.~\citep{voros2019limits} report a negative result of the
same kind: across five countries, friendship relations turn out to be a poor proxy
for status attributions, so one relation cannot be substituted for another. Both are
warnings that neighbour-based inference is conditional rather than automatic; our
contribution is to give one condition in closed form, for the case where the failure
is \emph{induced deliberately}.

Homophily is the assumption the inference rests on, and its measurement is a
developed literature here: Bojanowski and Corten~\citep{bojanowski2014measuring}
survey and axiomatise segregation and homophily indices, and
McMillan~\citep{mcmillan2022worth} shows that homophily measured on strong and weak
ties differs in magnitude, so that a single homophily figure for a population can be
misleading. Noel and Nyhan~\citep{noel2011unfriending} demonstrate the consequence for
inference, showing that homophily in tie \emph{retention} biases estimates of social
influence. Recent work in the journal continues in both directions: Sch\"u{\ss}ler and
Fuhse~\citep{schussler2025duality} argue that ties and attributes are mutually
constituting rather than one being read off the other, which is the conceptual
counterpart of our question; Espinosa-Rada et al.~\citep{espinosarada2026socioeconomic}
measure the socioeconomic composition of alters directly; and Tang et
al.~\citep{tang2026twomode} extend the autoregressive model to two-mode data with a
core--periphery degree structure, which is the setting our degree correction speaks to.
Where the network is only partially observed, Ajorlou et
al.~\citep{ajorlou2025dirichlet} give a sampling estimator for homophily itself.
Our results are the adversarial limit of this concern: we treat the degree
of homophily violation as a decision variable of one class rather than a property of
the population.

\subsection{Multiplexity, and why the analysis is per relation}
\label{sec:related-multiplex}

That different relations among the same actors have different structure is a founding
observation of this journal~\citep{lazega1999multiplexity}, and its methodological
consequence---that collapsing relations into a single adjacency destroys information,
while modelling them jointly quickly becomes intractable---is treated directly by
V\"or\"os and Snijders~\citep{voros2017cluster}. Snijders et
al.~\citep{snijders2013model} show the stronger version: structure apparently present
in an advice relation is substantially mediated by the friendship relation beneath
it. An et al.~\citep{an2026estimating} find the same phenomenon quantitatively, with
peer effects differing in magnitude between friends and classmates in a multilayer
setting.

Two recent contributions sharpen the point. Fluer et
al.~\citep{fluer2026survey} build multiplex models that carry interlayer
correlation explicitly rather than treating layers as independent, and Chandrasekhar et
al.~\citep{chandrasekhar2026multiplexing} show that multiplexity can have
opposite-signed consequences depending on the process running over it---inhibiting
low-reinforcement diffusion while enhancing high-reinforcement diffusion. That a
quantity can reverse sign across layer configurations is precisely what we find for the
usability of relational inference. Lakdawala et
al.~\citep{lakdawala2026bonds} meanwhile partition dyads into latent interaction
regimes, which is a complementary way of admitting that not all ties of a nominal type
are the same.

This matters for us empirically rather than only in principle.
Section~\ref{sec:measuring} finds camouflage differing by a factor approaching two
between relations recorded on the \emph{same} actors, which means our thresholds are
per-relation quantities and that a procedure pooling relations is averaging a usable
signal with an unusable one.

\subsection{Strategic tie formation}
\label{sec:related-strategic}

The premise that actors form ties instrumentally, and with an eye to how they will be
perceived, is established here but thinly populated. Burt et al.~\citep{burt2021network}
show experimentally that brokerage position causes attributions of leadership, and
report participants re-shaping an assigned network to suit their preferences. Shea et
al.~\citep{shea2019cheater} come closest to our mechanism: actors anticipating
unethical behaviour activate sparser networks that better conceal it, while those who
have already acted unethically activate denser ones for support---network structure
managed as a function of what the actor wants inferred about it. Barone and
Coscia~\citep{barone2018birds} estimate trustworthiness from mutually inconsistent
self-reports and find homophily in it, which is both an instance of inferring a
latent attribute that actors have an incentive to misreport and a demonstration that
the inference can work.

Nguyen~\citep{nguyen2026agency} adds the most recent instance, treating tie
formation as strategic behaviour conditional on structural position, and
Macdonald~\citep{macdonald2026wiretap} shows that covert networks carry
structural signatures distinct from ordinary collaboration---evidence that concealment
leaves traces in structure rather than erasing them.

What is absent, so far as we can determine, is the case where the tie portfolio is
\emph{engineered} to defeat the inference. We searched the journal's full record and
found no treatment of automated or inauthentic accounts; the nearest work concerns
misinformation susceptibility in personal networks~\citep{choi2026who}. We therefore
take our terminology from this journal---alter composition, autocorrelation,
exposure---and our substantive setting from the detection literature, rather than
importing the latter's vocabulary wholesale.

Within that detection literature a small strand does measure the structure of such
populations rather than only classifying them, and it is the strand closest to our
purpose: Ran et al.~\citep{ran2025identifying} characterise automated accounts by
motif frequencies across relation types, and Ariyarathne
et al.~\citep{ariyarathne2026behavior} treat measured behavioural change as the
discriminating signal in its own right. Two further results establish that camouflage is
an active and adaptive strategy rather than a fixed property: Pan et
al.~\citep{pan2026camera} document deliberate mimicry of benign users' content,
which breaks the assumption that structural and attribute cues agree, and Mukherjee
et al.~\citep{mukherjee2026optimal} show that detectors relying on aggregation can
be defeated by rewiring under realistic budget constraints. Our threshold gives the
condition under which such rewiring succeeds.

\subsection{Degree heterogeneity and the correction it forces}
\label{sec:related-degree}

Our edge-endpoint identity requires a degree correction, and the reason is a
long-standing concern here. Snijders~\citep{snijders1981degree} introduced degree
variance as an index of graph heterogeneity; Anderson et al.~\citep{anderson1999interaction}
showed that graph-level indices are systematically confounded by size and density; and
Krivitsky et al.~\citep{krivitsky2009representing} exhibit networks with identical,
highly skewed degree distributions but entirely different structure, concluding that
degree-based accounts cannot by themselves identify structure. Section~\ref{sec:measuring}
bears this out: the identity requires its degree factor wherever class mean degrees
differ, and holds exactly once it is included.

Kami\'nski et al.~\citep{kaminski2026relationships} extend the concern to
hypergraphs, finding an association between node degree and hyperedge size in all but a
small minority of $36$ empirical datasets, so degree is rarely separable from the
structure one wishes to measure.

The stochastic blockmodel we build on originates in this
journal~\citep{holland1983stochastic}. Degree correction \emph{within} blockmodels is
much less developed here; the one instance we found, Yu et
al.~\citep{yu2018detecting}, monitors node propensity in a dynamic degree-corrected
blockmodel in order to detect anomalous nodes, which is close to our setting in both
model and purpose.

\subsection{Analyses of neighbourhood averaging in the machine-learning literature}


The contextual stochastic block model~\citep{deshpande2018contextual} couples a
stochastic block model with Gaussian node attributes and has become the standard
setting for asking what message passing does to statistical separability. Baranwal
et al.~\citep{baranwal2021graph} established the central positive result: for a
two-class CSBM, one graph convolution improves the regime of linear separability by
a factor of roughly $1/\sqrt{D}$ in the expected degree $D$, and the resulting
classifier generalises out of distribution to graphs with different edge
probabilities.

That analysis assumes the generative process is homophilic---inter-class edges arise
from a fixed, non-adversarial probability---and establishes only the positive side.
Its follow-ups extend the positive results to multiple
layers~\citep{baranwal2023effects} and characterise the locally Bayes-optimal
architecture as one that \emph{interpolates} between an MLP in the low-graph-signal
regime and a convolution in the high-signal regime~\citep{baranwal2023optimality}.
Our Corollary~\ref{cor:critical} reduces to the separability gain
of~\citet{baranwal2021graph} when $\rho = 0$, which calibrates our model against this
baseline. That line has continued: Duranthon and
Zdeborov\'a~\citep{duranthon2025statistical} benchmark aggregation against the
information-theoretic optimum in the same generative model, quantifying how far short it
falls; Ma et al.~\citep{ma2025attention} show that attention---the closest
analogue to reweighting ties---helps only when structural noise dominates attribute
noise and hurts otherwise; and Baranwal~\citep{baranwal2026depth} obtains a sharp
dichotomy for the value of aggregation depth on sparse graphs, governed by a single
ratio. Our depth result in Section~\ref{sec:depth} is finite-$D$ and concerns the
crossover rather than the asymptotic regime, but it points the same way: depth is not
uniformly beneficial.

\paragraph{The crossover is known; our contribution is the imbalanced and
adversarial case.} We state the priority situation plainly, because it determines
what this paper does and does not claim. A closed-form threshold at which mean
aggregation becomes worse than the graph-free classifier was established by Ma
et al.~\citep[Thm.~2]{ma2022homophily} for the \emph{balanced} symmetric CSBM, and
Corollary~\ref{cor:critical} recovers their condition exactly at $c = 1$
(Remark~\ref{rem:ma} works the algebra through). Related crossover analyses have
since appeared: Luan et al.~\citep{luan2023when} compute the Bayes error of
graph-aware against graph-agnostic classifiers under a homophily parameter with
per-class variances and degrees, locating the intersection points directly and
criticising the fixed-classifier comparison used in~\citep{ma2022homophily}; Mao
et al.~\citep{mao2023demystifying} prove the per-node version, that a GCN beats an
MLP on homophilic nodes and loses on heterophilic ones \emph{within the same graph};
and the graph-attention analysis of Fountoulakis
et al.~\citep{fountoulakis2023graph} both identifies the same signal-deflation
factor and proves that attention---the closest analogue to soft edge
pruning---is lower-bounded by the better of convolution and the graph-free
classifier, while also showing (their Theorems~9 and~13) that no attention
mechanism separates intra- from inter-class edges in the hard regime.

Against this background our claim is narrower than ``we show when the graph is
harmful''. It is that the adversarial, class-imbalanced case is not a relabelling of
the balanced one: the two classes deflate asymmetrically, which moves the decision
boundary and invalidates a fixed-boundary comparison; bot prevalence enters the
threshold through $\eta = c\rho$ as a robustness parameter; and the mixture-variance
term omitted in~\citep{ma2022homophily} shifts the threshold by $11.5\%$. Read this way,
the closest antecedent to Theorem~\ref{thm:sgtr} is~\citep{fountoulakis2023graph},
which gives the qualitative guarantee for attention; we give a closed form for how
far a pruner of a given quality moves the threshold.

A parallel literature studies the limits of repeated aggregation. Over-smoothing
analyses~\citep{li2018deeper,oono2020graph} show node representations converge to
an uninformative subspace as depth grows, and spectral
accounts~\citep{nt2019revisiting} characterise graph convolution as a low-pass
filter. These results concern degradation with \emph{depth} under a fixed graph.
The effect we analyse is orthogonal: a single layer already suffers signal loss
when the \emph{topology itself} is adversarial. The two mechanisms compound, but
our critical ratio is present at $L = 1$.

\subsection{Heterophily in graph learning}

The finding that homophily-assuming GNNs degrade on heterophilic graphs motivated a
family of architectures that separate ego from neighbour representations, aggregate
over higher-order neighbourhoods, or combine low- and high-frequency
signals~\citep{zhu2020beyond,luan2022revisiting}. Ma et al.~\citep{ma2022homophily}
show in addition that homophily is not strictly necessary: GNNs can succeed on
heterophilic graphs when inter-class neighbourhood \emph{distributions} are
sufficiently distinguishable.

That observation matters for our setting, and we take a position on it. Under
adversarial camouflage the distinguishability Ma et al.\ rely on is precisely what
the adversary is optimising away: a bot that matches the human neighbourhood
distribution eliminates the signal their argument requires. The distinction between
benign heterophily---structured, and therefore exploitable---and adversarial
heterophily, engineered to be unexploitable, is in our view the conceptual point
separating bot detection from general heterophilic node classification. It is a
distinction about the \emph{provenance} of the inter-class edges, not about the
algebra of aggregation, which is why our threshold coincides with theirs in the
balanced case while the interpretation of the parameter differs.

\subsection{Robust aggregation and adversarial topology}
\label{sec:related-robust}

Camouflage is a contamination phenomenon, and two mature literatures have already
treated it as one.

The first replaces the mean with a robust estimator. In Byzantine-robust distributed
learning a fraction of workers report corrupted updates, and aggregation rules such
as Krum~\citep{blanchard2017machine} and coordinate-wise median or trimmed
mean~\citep{yin2018byzantine} bound the resulting error. The idea has been imported
into GNN neighbourhood aggregation directly: Geisler
et al.~\citep{geisler2020reliable} construct an aggregation with the maximal
breakdown point of $0.5$, so the bias stays bounded while fewer than half of a
node's edges are adversarial, and Chen et al.~\citep{chen2021understanding} attribute
the structural vulnerability of GCNs specifically to the non-robustness of the
weighted mean, analysed through the breakdown point.

We therefore avoid the term for our own threshold. A breakdown point asks when an
adversary can drive an estimate arbitrarily far under unbounded contamination;
$\rho^\star$ is instead a bias--variance crossover under \emph{bounded}
contamination---human features are not arbitrary---and marks where the $\sqrt{D}$
variance reduction ceases to pay for the $(1+c)\rho$ bias. The two quantities answer
different questions, and $\rho^\star$ is the weaker of the two.

The second literature attacks the topology. Nettack~\citep{zugner2018adversarial} and
Metattack~\citep{zugner2019adversarial} insert or remove edges to flip predictions,
and defences such as GNNGuard~\citep{zhang2020gnnguard} and
Pro-GNN~\citep{jin2020graph} filter them, largely heuristically. Two results in this
line bear directly on ours. Zhu et al.~\citep{zhu2022how} prove the bridge our model
assumes---that impactful structural attacks on homophilous data necessarily reduce
homophily---and derive the ego/neighbour-separation design principle from it.
Mujkanovic et al.~\citep{mujkanovic2022are} show empirically that under adaptive
attack, defended GNNs can underperform an MLP that ignores the topology outright;
our Corollary~\ref{cor:critical} can be read as a closed-form account of when that
should be expected. Certified-robustness work~\citep{bojchevski2019certifiable} is
orthogonal: it certifies that a prediction is unchanged within a perturbation
budget, never that the graph has become net-harmful beyond one. Gosch
et al.~\citep{gosch2023revisiting} come closest in combining a CSBM, an adversary,
and a Bayes-optimal yardstick, but use it to show GNNs are \emph{over}-robust,
insensitive past the point of genuine semantic change.

\subsection{Camouflage-resistant detectors}

CARE-GNN~\citep{dou2020caregnn} introduced the camouflage framing to GNN-based fraud
detection, using a label-aware similarity measure and a reinforcement-learning
neighbour selector to filter uninformative edges. Bot-detection work has since
adopted the pattern: BotHP~\citep{he2025bothp} separates ego and neighbour encoders
to mitigate interaction camouflage, and XHBot~\citep{dang2026xhbot} combines
spectral edge pruning, tri-channel aggregation, and contrastive disentanglement.

These methods are validated by ablation. An ablation establishes that removing a
component costs accuracy on a particular benchmark; it does not establish the
\emph{range of conditions} under which the component helps, nor whether a
multi-channel design is more expressive or merely better parameterised. The present
paper supplies exactly those missing statements for the three mechanisms XHBot
employs. We analyse XHBot's components because they instantiate the dominant design
pattern, not because the analysis is specific to that system: the edge-scoring
result applies to any feature-discrepancy pruning rule, and the expressivity result
to any aggregator that adds a permutation-invariant non-mean channel.



\section{The Camouflage-CSBM}
\label{sec:model}

We work in a two-class contextual stochastic block model augmented with an
adversarial parameter controlling how aggressively bots attach to humans.

\begin{definition}[Camouflage-CSBM]
\label{def:ccsbm}
Let $\mathcal{V}$ be a set of $N$ nodes partitioned into bots $\mathcal{B}$ and
humans $\mathcal{H}$, with bot prevalence $\pi_B = |\mathcal{B}|/N \in (0,1)$ and
bot-to-human ratio $c = \pi_B/(1-\pi_B)$. Each node $v$ carries a feature vector
\begin{equation}
    \mathbf{x}_v \sim \mathcal{N}(\boldsymbol{\mu}_{y_v},\, \sigma^2 I_d),
    \qquad y_v \in \{\mathrm{B}, \mathrm{H}\},
\end{equation}
and we write $\Delta = \boldsymbol{\mu}_{\mathrm{B}} - \boldsymbol{\mu}_{\mathrm{H}}$
for the class-mean gap. Every node has neighbourhood size $D$. Neighbours are drawn
independently: a bot's neighbour is human with probability $\rho$ (the
\emph{camouflage ratio}) and a bot otherwise; a human's neighbour is a bot with
probability $\eta$ and a human otherwise.
\end{definition}

The two mixing parameters are not free. Counting bot--human edge endpoints from
both sides gives $|\mathcal{B}| D \rho = |\mathcal{H}| D \eta$, hence the
\emph{edge-balance identity}
\begin{equation}
    \eta(\rho) \;=\; \frac{\pi_B}{1-\pi_B}\,\rho \;=\; c\,\rho .
    \label{eq:balance}
\end{equation}
Camouflage is therefore doubly harmful: it not only pollutes bot neighbourhoods but
necessarily pollutes human neighbourhoods as well, at a rate set by the bot
prevalence. Every result below accounts for this coupling; treating $\eta$ as
independent of $\rho$ overstates a detector's robustness.

\begin{assumption}[Regularity used by the depth analysis]
\label{as:tree}
Definition~\ref{def:ccsbm} samples each actor's alters independently, which suffices
for the single-layer results. Section~\ref{sec:depth} additionally requires that the
tie relation be \emph{symmetric}---so that $v \in \mathcal{N}(u)$ whenever
$u \in \mathcal{N}(v)$, which is what makes~\eqref{eq:backtrack} hold---and that the
neighbourhood be locally tree-like with branching $D$, so that the number of actors at
distance $j$ is $N_j = D(D-1)^{j-1}$. Both are standard for sparse networks in which
short cycles are rare, and neither is needed before Section~\ref{sec:depth}. We also
take $\rho \le 1/c$, so that $\eta \le 1$, which holds automatically whenever the
strategic class is the minority.
\end{assumption}

\begin{remark}[Modelling choices]
Independent neighbour sampling at fixed size $D$ makes the analysis exact and is
the discrete analogue of the degree-concentration assumptions used
in~\citep{baranwal2021graph}. Isotropic within-class covariance and a common $D$ are
simplifications; Section~\ref{sec:limitations} discusses what changes without them.
We emphasise that $\rho$ is an \emph{adversarial} parameter chosen by the attacker,
not a fixed property of the generative process---this is the substantive
difference from the standard CSBM setting.
\end{remark}

\paragraph{Aggregation operator.}
We analyse one round of mean aggregation, the canonical homophilic operator,
\begin{equation}
    \mathbf{h}_v \;=\; \frac{1}{D}\sum_{u \in \mathcal{N}(v)} \mathbf{x}_u ,
    \label{eq:meanagg}
\end{equation}
which isolates the effect of neighbourhood averaging from the confounds of learned
weights and nonlinearity. Because all the quantities we need are governed by the
projection onto the discriminant direction, define
$\mathbf{w} = \Delta/\|\Delta\|$ and $t_v = \mathbf{w}^\top \mathbf{h}_v$.

\paragraph{Baseline for comparison.}
The natural point of reference is the graph-free classifier, which uses
$\mathbf{x}_v$ directly. For two isotropic Gaussians with equal priors its Bayes
error is
\begin{equation}
    \varepsilon_{\mathrm{raw}} \;=\; \Phi\!\left(-\frac{\|\Delta\|}{2\sigma}\right),
    \label{eq:rawerr}
\end{equation}
with $\Phi$ the standard normal CDF. The question driving Section~\ref{sec:theory}
is when aggregation beats~\eqref{eq:rawerr} and when it does not.



\section{Theory}
\label{sec:theory}

\subsection{Camouflage deflates the signal}

We first quantify what aggregation does to the class signal.

\begin{proposition}[Signal deflation]
\label{prop:deflation}
Under Definition~\ref{def:ccsbm} with mean aggregation~\eqref{eq:meanagg}, the
expected post-aggregation separation along the discriminant direction is
\begin{equation}
    \delta(\rho) \;=\; \mathbb{E}[t_v \mid y_v = \mathrm{B}]
                     - \mathbb{E}[t_u \mid y_u = \mathrm{H}]
                \;=\; \bigl(1 - \rho - \eta(\rho)\bigr)\,\|\Delta\| ,
    \label{eq:deflation}
\end{equation}
and with the edge-balance identity~\eqref{eq:balance},
$\delta(\rho) = \bigl(1 - (1+c)\rho\bigr)\|\Delta\|$.
\end{proposition}

\begin{proof}
Centre the projection so that $\mathbf{w}^\top\boldsymbol{\mu}_{\mathrm{B}} =
+\|\Delta\|/2$ and $\mathbf{w}^\top\boldsymbol{\mu}_{\mathrm{H}} = -\|\Delta\|/2$.
For a bot $v$, each of its $D$ neighbours is human with probability $\rho$, so by
linearity
$\mathbb{E}[t_v] = (1-\rho)\tfrac{\|\Delta\|}{2} + \rho\bigl(-\tfrac{\|\Delta\|}{2}\bigr)
= \tfrac{\|\Delta\|}{2}(1-2\rho)$.
Symmetrically $\mathbb{E}[t_u] = -\tfrac{\|\Delta\|}{2}(1-2\eta)$ for a human $u$.
Subtracting gives $\tfrac{\|\Delta\|}{2}\bigl[(1-2\rho)+(1-2\eta)\bigr] =
(1-\rho-\eta)\|\Delta\|$. Substituting~\eqref{eq:balance} yields the second form.
\end{proof}

Two consequences are worth stating. The signal vanishes entirely at
$\rho_{\mathrm{null}} = 1/(1+c)$, at which point the two classes have identical
post-aggregation means and \emph{no} classifier acting on aggregated features
alone can do better than chance. And the deflation is linear in $\rho$ with slope
$(1+c)\|\Delta\|$, so higher bot prevalence makes the detector more fragile, not
less---a counterintuitive point, since prevalence is usually treated as a
class-imbalance nuisance rather than a robustness parameter.

The noise term requires more care than the signal term.

\begin{theorem}[Post-aggregation error]
\label{thm:error}
Under Definition~\ref{def:ccsbm}, the projected statistics are asymptotically
Gaussian with separation $\delta(\rho)$ from~\eqref{eq:deflation} and per-class
variances
\begin{equation}
    s_{\mathrm{B}}^2 = \frac{\sigma^2 + \rho(1-\rho)\|\Delta\|^2}{D},
    \qquad
    s_{\mathrm{H}}^2 = \frac{\sigma^2 + \eta(1-\eta)\|\Delta\|^2}{D}.
    \label{eq:variances}
\end{equation}
Defining the post-aggregation signal-to-noise ratio
\begin{equation}
    \mathrm{SNR}(\rho) \;=\; \frac{|\delta(\rho)|}{s_{\mathrm{B}} + s_{\mathrm{H}}},
    \label{eq:snr}
\end{equation}
the balanced error of the variance-equalised linear rule is
$\Phi(-\mathrm{SNR}(\rho))$. This is an upper bound on the Bayes error, tight when
$s_{\mathrm{B}} = s_{\mathrm{H}}$.
\end{theorem}

\begin{remark}[What this quantity is, and is not]
\label{rem:not-bayes}
Three clarifications, since the distinction matters for how the comparisons below
should be read. (i) When $s_{\mathrm{B}} \neq s_{\mathrm{H}}$ the likelihood ratio of
two unequal-variance Gaussians is quadratic, so the Bayes rule is an interval and
$\Phi(-\mathrm{SNR})$ is the error of the equal-error-rate \emph{threshold} rule, an
upper bound. That bound is loose in one direction, but at finite $D$ it is not the
dominant approximation: the central limit theorem is. The aggregate is exactly a
mixture of $D+1$ Gaussians, and $\Phi(-\mathrm{SNR})$ \emph{under}-states its error
by $41\%$ relative at $\rho = 0.1$, $18\%$ at $\rho = 0.2$ and $6\%$ at
$\rho = 0.3$ (Table~\ref{tab:deflation}), while the attribute-only baseline is exact
because raw attributes are exactly Gaussian. Comparisons therefore \emph{overstate}
how much aggregation helps, and are optimistic rather than conservative. The
crossover itself is barely affected---exact $0.4642$ against $0.4644$---but the
direction should not be misreported.
(ii) The criterion is \emph{balanced} error, weighting the two classes equally, not
the prior-weighted risk---a deliberate choice in a paper about class imbalance,
where a prior-weighted criterion rewards predicting the majority class and would
obscure the effect being measured. At $D=16$, $\pi_B=0.3$ the two differ little in
practice ($\rho^\star = 0.4644$ balanced against $0.4694$ prior-weighted), but the
choice should be explicit. (iii) The absolute value in~\eqref{eq:snr} is necessary:
without it the expression exceeds $1/2$ once $\rho$ passes the point where
$\delta$ changes sign, which no error probability may do. With it,
$\Phi(-\mathrm{SNR})$ correctly records that strongly anti-correlated neighbourhoods
are informative again---the effect made explicit in~\eqref{eq:window}.
\end{remark}

\begin{proof}
Condition on a single neighbour of a bot. Its projected feature equals
$\pm\|\Delta\|/2$ according to a Bernoulli$(\rho)$ draw, plus independent
$\mathcal{N}(0,\sigma^2)$ noise. By the law of total variance, the mean component
contributes $\|\Delta\|^2\rho(1-\rho)$ and the noise contributes $\sigma^2$, so each
neighbour has variance $\sigma^2 + \rho(1-\rho)\|\Delta\|^2$. The $D$ neighbours are
independent, so averaging divides the variance by $D$, giving $s_{\mathrm{B}}^2$;
the human case is identical with $\eta$ in place of $\rho$. Asymptotic normality of
$t_v$ follows from the central limit theorem as $D$ grows. For two equal-prior
Gaussians with common standard deviation $s$ and mean gap $\delta$, the
likelihood-ratio test thresholds at the midpoint and attains error
$\Phi(-\delta/2s)$, which is~\eqref{eq:snr} when $s_{\mathrm{B}} = s_{\mathrm{H}} = s$.
\end{proof}

\begin{remark}
The mixture term $\rho(1-\rho)\|\Delta\|^2$ is easy to miss: a calculation that
models only the additive noise obtains $\sigma^2/D$ and therefore
\emph{under}-estimates the post-aggregation error, most severely at intermediate
$\rho$ where the mixture variance peaks. Section~\ref{sec:validation} confirms
the corrected form matches simulation to four decimal places.
\end{remark}

\subsection{The critical camouflage ratio}

We can now answer the deployment question directly.

\begin{corollary}[Critical camouflage ratio]
\label{cor:critical}
Mean aggregation yields a lower Bayes error than the graph-free
classifier~\eqref{eq:rawerr} if and only if $\mathrm{SNR}(\rho) > \|\Delta\|/(2\sigma)$.
Writing $\rho^\star$ for the value of $\rho$ at which equality holds, aggregation
helps precisely when $\rho < \rho^\star$. In the noise-dominated regime
$\sigma^2 \gg \rho(1-\rho)\|\Delta\|^2$ this takes the closed form
\begin{equation}
    \rho^\star \;=\; \frac{1 - D^{-1/2}}{1+c}.
    \label{eq:critical}
\end{equation}
\end{corollary}

\begin{proof}
The ``if and only if'' is immediate from monotonicity of $\Phi$. In the
noise-dominated regime the mixture terms in~\eqref{eq:variances} are negligible, so
$s_{\mathrm{B}} \approx s_{\mathrm{H}} \approx \sigma/\sqrt{D}$ and
$\mathrm{SNR}(\rho) \approx (1-\rho-\eta)\|\Delta\|\sqrt{D}/(2\sigma)$. Setting this
equal to $\|\Delta\|/(2\sigma)$ gives $(1-\rho-\eta)\sqrt{D} = 1$; substituting
$\eta = c\rho$ and solving for $\rho$ yields~\eqref{eq:critical}.
\end{proof}

Equation~\eqref{eq:critical} is the paper's most directly actionable result, and it
says three things. First, at $\rho = 0$ the condition reduces to
$\mathrm{SNR}_{\mathrm{agg}}/\mathrm{SNR}_{\mathrm{raw}} = \sqrt{D}$: aggregation
improves separability by $\sqrt{D}$, recovering the $1/\sqrt{D}$ gain established
by Baranwal et al.~\citep{baranwal2021graph} and confirming our model is calibrated
against the known homophilic baseline. Second, $\rho^\star$ increases with $D$, so
detectors operating on densely connected accounts tolerate more camouflage---the
graph is a genuine asset for hub accounts long after it has become a liability for
sparse ones. Third, $\rho^\star$ decreases in $c$: as bots become more prevalent,
the camouflage budget they need in order to break the detector shrinks.

\begin{remark}[Relation to Ma et al.]
\label{rem:ma}
Equation~\eqref{eq:critical} is not new in the balanced case, and we want to be
explicit about what is and is not ours. For a symmetric CSBM with intra- and
inter-class edge probabilities $p$ and $q$, Ma et al.~\citep[Thm.~2]{ma2022homophily}
prove that a node is better classified after aggregation than before whenever
$\deg(i) > (p+q)^2/(p-q)^2$. Their inter-class neighbour fraction is
$\rho = q/(p+q)$, so $(p-q)/(p+q) = 1-2\rho$ and their condition rearranges to
$\sqrt{D}\,(1-2\rho) > 1$, i.e.\ $\rho < (1 - D^{-1/2})/2$. This is
exactly~\eqref{eq:critical} at $c = 1$. A closed-form threshold for when mean
aggregation is worse than discarding the graph was therefore established at
ICLR~2022, and our corollary is its extension to $c \neq 1$.

What the extension adds is not merely a symbol. Three things change when the classes
are imbalanced. (i) The two classes deflate by \emph{different} amounts, $\rho$ for
bots and $\eta = c\rho$ for humans, so the post-aggregation class midpoint
\emph{moves}; a comparison against a decision boundary held fixed at the raw-feature
midpoint---as in~\citep{ma2022homophily}---is then no longer the right one, which is
precisely the criticism Luan et al.~\citep{luan2023when} raise against that
analysis. Our Theorem~\ref{thm:error} compares Bayes errors instead, and is
insensitive to this. (ii) The edge-balance identity $\eta = c\rho$
in~\eqref{eq:balance} is what couples the two, and it is what makes bot
prevalence a robustness parameter rather than a training-time nuisance. (iii) The
mixture-variance term $\rho(1-\rho)\|\Delta\|^2$ is absent from the covariance
$I/\sqrt{\deg(i)}$ reported in~\citep{ma2022homophily}; including it moves
$\rho^\star$ from $0.5250$ to $0.4644$ in the configuration of
Remark~\ref{rem:practical}, an $11.5\%$ correction to the threshold.
\end{remark}

\begin{remark}[Practical reading]
\label{rem:practical}
For a detector with typical $D = 16$ and bot prevalence $\pi_B = 0.3$
($c \approx 0.43$), \eqref{eq:critical} gives $\rho^\star \approx 0.5250$ in the
noise-dominated regime and $0.4644$ once the mixture variance is included---an
$11.5\%$ correction. A bot that
directs roughly half its edges at genuine users has therefore neutralised the graph
entirely. This is not an extreme adversary; it is a modest one, which is why
architectural mitigation is necessary rather than optional.
\end{remark}

\subsection{Spectral edge pruning provably extends the operating regime}
\label{sec:sgtr-theory}

The spectral tie-scoring rule of \citet{dang2026xhbot} scores each edge by a learned function of the feature
discrepancy $(\mathbf{x}_u - \mathbf{x}_v)^2$, motivated by the per-edge
decomposition of the Dirichlet energy $\mathbf{f}^\top L \mathbf{f} =
\sum_{(u,v)\in\mathcal{E}} (f_u-f_v)^2$. The following lemma shows this quantity is
the right statistic: it separates camouflage from genuine edges by exactly the
squared class gap, independently of the feature dimension.

\begin{lemma}[Energy separability]
\label{lem:energy}
Under Definition~\ref{def:ccsbm}, for a within-class edge $(u,v)$ and a
cross-class (camouflage) edge $(u',v')$,
\begin{equation}
    \mathbb{E}\bigl\|\mathbf{x}_u - \mathbf{x}_v\bigr\|^2 = 2d\sigma^2,
    \qquad
    \mathbb{E}\bigl\|\mathbf{x}_{u'} - \mathbf{x}_{v'}\bigr\|^2 = \|\Delta\|^2 + 2d\sigma^2,
\end{equation}
so the expected energy gap is exactly $\|\Delta\|^2$.
\end{lemma}

\begin{proof}
For a within-class edge both endpoints share a mean, so
$\mathbf{x}_u - \mathbf{x}_v \sim \mathcal{N}(0, 2\sigma^2 I_d)$ and the expected
squared norm is $2d\sigma^2$. For a cross-class edge
$\mathbf{x}_{u'} - \mathbf{x}_{v'} \sim \mathcal{N}(\Delta, 2\sigma^2 I_d)$, and
$\mathbb{E}\|Z\|^2 = \|\mathbb{E}Z\|^2 + \mathrm{tr}\,\mathrm{Cov}(Z)$ gives
$\|\Delta\|^2 + 2d\sigma^2$.
\end{proof}

The gap is a \emph{location} shift that does not grow with $d$, whereas the common
term $2d\sigma^2$ does; the score is therefore informative but increasingly noisy in
high dimension, which is precisely why XHBot learns a projection rather than using
the raw quadratic form. We now quantify the benefit of acting on this score.

\begin{theorem}[spectral tie scoring raises the tolerated camouflage]
\label{thm:sgtr}
Suppose soft pruning multiplies camouflage edge weights by $(1-\bar e_c)$ and
within-class edge weights by $(1-\bar e_w)$ in expectation, with
$\bar e_c > \bar e_w$. Then weighted aggregation behaves as unweighted aggregation
at the \emph{effective} camouflage ratio
\begin{equation}
    \rho_{\mathrm{eff}}(\rho)
    \;=\; \frac{\rho(1-\bar e_c)}{\rho(1-\bar e_c) + (1-\rho)(1-\bar e_w)}
    \;<\; \rho .
    \label{eq:rhoeff}
\end{equation}
Consequently, with $\bar e_w = 0$, the detector tolerates true camouflage up to
\begin{equation}
    \rho^\star_{\mathrm{prune}}
    \;=\; \frac{\rho^\star}{\rho^\star + (1-\bar e_c)(1-\rho^\star)}
    \;\geq\; \rho^\star,
    \label{eq:rhostarsgtr}
\end{equation}
with equality if and only if $\bar e_c = 0$.
\end{theorem}

\noindent
Equation~\eqref{eq:rhostarsgtr} accounts only for the change in neighbourhood
\emph{composition}. It therefore appears to imply
$\rho^\star_{\mathrm{prune}} \to 1$ as $\bar e_c \to 1$: that a sufficiently good
pruner tolerates any camouflage whatsoever. That conclusion is false, and
Section~\ref{sec:achievable} shows why. Down-weighting edges non-uniformly also
reduces the \emph{effective degree}, and hence forfeits part of the $\sqrt{D}$
variance reduction that made aggregation worthwhile in the first place. Once that
cost is included the naive reading fails, though not in the direction it first
appears: Proposition~\ref{cor:ceiling} shows there is no hard ceiling at all, because
a perfect pruner leaves the aggregate unbiased. What changes is the \emph{rate} at
which the advantage thins out, and the fact that the operating point must be chosen on
a detection trade-off rather than by pushing $\bar e_c \to 1$.

\begin{proof}
Weighted mean aggregation computes $\sum_e w_e \mathbf{x}_{u_e} / \sum_e w_e$. In
expectation a fraction $\rho$ of a bot's edges are camouflage edges carrying weight
$(1-\bar e_c)$ and $(1-\rho)$ are within-class edges carrying $(1-\bar e_w)$;
normalising gives the weight share~\eqref{eq:rhoeff} on human neighbours, which is
exactly the role $\rho$ plays in Proposition~\ref{prop:deflation}. The inequality
$\rho_{\mathrm{eff}} < \rho$ follows since $\bar e_c > \bar e_w$ makes the numerator
shrink faster than the denominator. For~\eqref{eq:rhostarsgtr}, set
$\rho_{\mathrm{eff}}(\rho) = \rho^\star$ with $\bar e_w = 0$ and solve: writing
$\beta = 1-\bar e_c$, the equation $\rho\beta = \rho^\star(\rho\beta + 1 - \rho)$
rearranges to $\rho[\beta(1-\rho^\star) + \rho^\star] = \rho^\star$, giving the
stated form. Monotonicity in $\bar e_c$ and the limits are immediate.
\end{proof}

Theorem~\ref{thm:sgtr} converts an ablation delta into a statement about the
operating regime: a pruning module that attenuates camouflage edges by $75\%$ on
average lifts the tolerated camouflage from $0.53$ to $0.82$ at $D=16$,
$\pi_B = 0.3$. It also clarifies what pruning cannot do. Because
$\rho_{\mathrm{eff}} \to 1$ as $\rho \to 1$ for any fixed $\bar e_c < 1$, no
pruning rule with imperfect discrimination defeats a maximally camouflaged
adversary; the spectral rule buys operating range, not immunity.

\begin{remark}[What is new here, and what is not]
\label{rem:sgtr-priority}
The \emph{mechanism} in~\eqref{eq:rhoeff}---that edge reweighting acts by rescaling
the effective inter-class edge share, so that weights enter only through the products
of weight and edge rate---is not new. Weiss~\citep[Thm.~2]{weiss2026cost} proves the
corresponding statement for cost-sensitive aggregation with within- and
cross-class weights $w_+ > w_-$, and shows the class-mean direction is preserved iff
$w_+/w_- > q/p$; our~\eqref{eq:rhoeff} is that identity with $w_+ = 1-\bar e_w$ and
$w_- = 1-\bar e_c$. The observation that only a perfect discriminator tolerates
arbitrary heterophily is the contrapositive of the same condition. The related
guarantee for graph attention in~\citep{fountoulakis2023graph} is qualitative but
similar in spirit.

What~\eqref{eq:rhostarsgtr} adds is the \emph{inversion}: composing that rescaling
with the finite-$D$ threshold of Corollary~\ref{cor:critical} to obtain a bound on the
\emph{adversary's} budget as an explicit function of pruner quality $\bar e_c$. This
is a short derivation, and we present it as a corollary of two known ingredients
rather than as an independent discovery. Its value is operational: $\bar e_c$ is
measurable for a deployed pruner, so~\eqref{eq:rhostarsgtr} turns a component
ablation into a quantitative operating-range claim. We note also
that~\citep{weiss2026cost} works in the dense regime where degree effects are absorbed
into $O(1/n)$, which removes the $D^{-1/2}$ term that our threshold depends on; the
finite-$D$ setting is therefore necessary for a statement of this form.
\end{remark}

\subsection{What pruning quality is achievable, and what it costs}
\label{sec:achievable}

Theorem~\ref{thm:sgtr} takes the pruner's quality $\bar e_c$ as given. Two things
are then missing: whether a given $\bar e_c$ is attainable at all by a
feature-based score, and what is given up in exchange. We supply both, which
removes the free parameter and corrects the limit noted above.

\paragraph{Only the retained-weight ratio matters.}
Write $a = 1-\bar e_c$ and $b = 1-\bar e_w$ for the surviving weight on camouflage
and within-class edges. Weighted mean aggregation is invariant to a common rescaling
of all weights, so $a$ and $b$ enter only through the ratio
\begin{equation}
    \theta \;=\; a/b \;\in\; (0,1] ,
    \label{eq:theta}
\end{equation}
which we call the \emph{retained-weight ratio}; $\theta = 1$ is no pruning and
$\theta = 0$ is perfect pruning. This is the same quantity that governs the
cost-sensitive condition $w_+/w_- > q/p$ of~\citep{weiss2026cost}. In terms of
$\theta$, the effective camouflage ratio of Theorem~\ref{thm:sgtr} is
$\rho_{\mathrm{eff}} = \rho\theta/(\rho\theta + 1-\rho)$.

\begin{proposition}[Pruning forfeits effective degree]
\label{prop:deff}
A weighted mean of $D$ unit-variance terms has the variance of an unweighted mean
of
\begin{equation}
    D_{\mathrm{eff}}
    \;=\; D\,\frac{\bigl[(1-\rho) + \rho\theta\bigr]^2}
                   {(1-\rho) + \rho\theta^{2}}
    \;\leq\; D
    \label{eq:deff}
\end{equation}
terms, with equality if and only if $\theta = 1$.

Pruning acts on \emph{both} ends of a camouflage edge, so the human side is
rescaled too and its cross-class share becomes
$\eta_{\mathrm{eff}} = \eta\theta/(\eta\theta + 1-\eta)$ with $\eta = c\rho$. Note
that $\eta_{\mathrm{eff}} \neq c\,\rho_{\mathrm{eff}}$ unless $c=1$ or $\theta=1$:
the edge-balance identity holds for the \emph{raw} ratios and is not preserved by
pruning. Writing
$\lambda_{\mathrm{eff}} = 1 - \rho_{\mathrm{eff}} - \eta_{\mathrm{eff}}$ and
$D^{B}_{\mathrm{eff}}, D^{H}_{\mathrm{eff}}$ for~\eqref{eq:deff} evaluated at $\rho$
and at $\eta$ respectively, pruned aggregation beats the graph-free classifier if and
only if
\begin{equation}
    \frac{|\lambda_{\mathrm{eff}}|}
         {\sqrt{\tfrac12\bigl(1/D^{B}_{\mathrm{eff}} + 1/D^{H}_{\mathrm{eff}}\bigr)}}
    \;>\; 1 .
    \label{eq:pruned-condition}
\end{equation}
\end{proposition}

\begin{proof}
The variance of $\sum_e w_e z_e / \sum_e w_e$ for independent unit-variance $z_e$ is
$\sum_e w_e^2/(\sum_e w_e)^2$, so the variance-equivalent count is
$(\sum_e w_e)^2/\sum_e w_e^2$. Substituting $\rho D$ edges of weight $a$ and
$(1-\rho)D$ of weight $b$, and cancelling $b^2$, gives~\eqref{eq:deff}; the
inequality is Cauchy--Schwarz, with equality iff the weights are constant.
Combining with Proposition~\ref{prop:deflation} applied at $\rho_{\mathrm{eff}}$
yields~\eqref{eq:pruned-condition}.
\end{proof}

This resolves a degeneracy that~\eqref{eq:rhostarsgtr} alone does not. Uniform
attenuation ($\theta = 1$, any common scale) leaves both $\rho_{\mathrm{eff}}$ and
$D_{\mathrm{eff}}$ unchanged and so achieves exactly nothing, as it must; and
aggressive pruning is now penalised rather than free.

\begin{proposition}[A perfect pruner never loses, but its benefit thins out]
\label{cor:ceiling}
Under perfect pruning every retained neighbour shares the node's label, so the
retained count is $K \sim \mathrm{Bin}(D, 1-\rho)$ and the aggregate is
\emph{unbiased} with variance $\sigma^2/K$. Consequently, for every $\rho<1$,
pruned aggregation is never worse than the graph-free classifier: it is strictly
better on the event $\{K\geq2\}$, and ties when $K\in\{0,1\}$. The fraction of \emph{strategic actors}
that strictly benefit is
\begin{equation}
    \Pr[K \geq 2] \;=\; 1 - \rho^{D} - D(1-\rho)\rho^{D-1},
    \label{eq:benefit}
\end{equation}
and the variance-equivalent degree is $1/\mathbb{E}[1/K \mid K\geq1]$, not
$\mathbb{E}[K]$.
\end{proposition}

\begin{proof}
With all camouflage edges removed the surviving neighbours are same-class, so their
mean has the node's own class mean as its expectation: the bias is zero regardless of
$\rho$. Conditional on $K=k\geq1$ the variance is $\sigma^2/k \leq \sigma^2$, with
strict inequality iff $k\geq2$; nodes with $K=0$ have no neighbourhood and fall back
on their own features, tying exactly. Since $K$ is Binomial, $\Pr[K\leq1]$ is the sum
of its first two terms, giving~\eqref{eq:benefit}. That the variance-equivalent
degree is the harmonic rather than arithmetic mean follows because the variance of a
mean of $K$ terms is $\sigma^2\mathbb{E}[1/K]$, and
$\mathbb{E}[1/K] \geq 1/\mathbb{E}[K]$ by Jensen.
\end{proof}

\begin{remark}[A tempting but incorrect ceiling]
\label{rem:no-ceiling}
Substituting the \emph{mean} retained count $D(1-\rho)$
into~\eqref{eq:deff} yields $D_{\mathrm{eff}} = D(1-\rho)$ and hence the clean-looking
conclusion that even perfect pruning requires $\rho < 1-1/D$. That is wrong, and
instructively so. The substitution replaces a random $K$ by its mean in a quantity
that depends on $1/K$, and it is applied in precisely the regime where the
replacement is least defensible: at $\rho = 1-1/D$ and $D=16$ we have
$\mathbb{E}[K] = 1.00$ exactly, yet $1/\mathbb{E}[1/K \mid K\geq1] = 1.29$, and $36\%$ of
nodes have $K=0$. Proposition~\ref{cor:ceiling} shows the qualitative conclusion
inverts: there is no hard ceiling. What does happen as $\rho \to 1$ is that the
\emph{proportion} of nodes for which the graph carries any information at all decays;
a useful summary is the point at which half of the strategic actors still benefit,
which is $\rho = 0.897$ at $D=16$ and $0.974$ at $D=64$.
\end{remark}

\paragraph{Achievable $\theta$.}
It remains to bound $\theta$ from below. The spectral rule scores an edge by the feature
discrepancy $s_{uv} = \|\mathbf{x}_u - \mathbf{x}_v\|^2$, whose law under the CSBM
is exactly
\begin{align}
    s &\sim 2\sigma^2\chi^2_{d}
    \quad\text{(within class)}, \nonumber\\
    s &\sim 2\sigma^2\chi^2_{d}\!\left(\tfrac{r^2}{2}\right)
    \quad\text{(camouflage)},
    \label{eq:scorelaw}
\end{align}
with $r = \|\Delta\|/\sigma$ and $\chi^2_d(\lambda)$ the noncentral chi-square. Any
threshold rule therefore has an operating point $(\bar e_w, \bar e_c)$ on the ROC
of~\eqref{eq:scorelaw}, and $\theta$ is the ratio of its complements. In the
large-$d$ Gaussian approximation the discriminability index is
\begin{equation}
    \kappa \;=\; \frac{\|\Delta\|^2}{2\sigma^2\sqrt{2d}}
            \;=\; \frac{r^2}{2\sqrt{2d}} ,
    \label{eq:kappa}
\end{equation}
so $\bar e_c \approx \Phi(\Phi^{-1}(\bar e_w) + \kappa)$.

Two consequences follow. First, $\bar e_w \to 0$ forces $\bar e_c \to 0$: there is no
threshold that removes camouflage edges without also removing genuine ones, so the
idealisation $\bar e_w = 0$ used in~\eqref{eq:rhostarsgtr} is not attainable and the
joint form is required. Second, because $\kappa$ scales as $d^{-1/2}$, the raw
discrepancy score \emph{degrades} as the feature dimension grows---the two score
distributions concentrate and overlap. This is the formal content of the remark that
motivates learning a projection: mapping to $k \ll d$ dimensions while preserving
$\|\Delta\|$ multiplies $\kappa$ by $\sqrt{d/k}$.

\paragraph{The soft-weight and hard-threshold models must not be mixed.}
Equation~\eqref{eq:deff} assigns every camouflage edge the same weight $a$ and every
genuine edge the same $b$, so it is invariant to a common rescaling and depends only
on $\theta$. A \emph{threshold} rule is different: it deletes edges outright, so the
neighbourhood shrinks in absolute terms and the effective degree is
\begin{equation}
    D^{\mathrm{hard}}_{\mathrm{eff}}
    = D\bigl[(1-\rho)(1-\bar e_w) + \rho\,(1-\bar e_c)\bigr] ,
    \label{eq:deff-hard}
\end{equation}
which collapses as $\bar e_w \to 1$. Optimising the operating point under the
scale-invariant model~\eqref{eq:deff} while reading $(\bar e_w, \bar e_c)$ off a
threshold ROC is therefore inconsistent, and the inconsistency is not benign: it
rewards discarding almost every genuine edge at no modelled cost, so the apparent
optimum runs to $\bar e_w \to 1$ and no interior solution exists. We therefore
optimise~\eqref{eq:deff-hard} together with the ROC of~\eqref{eq:scorelaw}, which has
a genuine interior optimum and is stable under enlargement of the search domain.
The resulting tolerance depends only on $(D, c, r, d)$, all measurable.
Section~\ref{sec:validation} reports the values, which are considerably less
flattering to raw-discrepancy pruning than the inconsistent pairing suggests.

\subsection{Depth}
\label{sec:depth}

Everything so far concerns a single layer. Since deployed detectors stack two or
more, we record how the picture changes. The analysis requires care, because the
obvious recursion---each layer multiplying the separation by $\lambda$---is
\emph{not} what mean aggregation does.

\paragraph{Mean aggregation backtracks.}
On a $D$-regular tree, two rounds of neighbourhood averaging give
\begin{equation}
    \mathbf{h}^{(2)}_v
    = \frac{1}{D^2}\sum_{u\in\mathcal{N}(v)}\sum_{w\in\mathcal{N}(u)}\mathbf{x}_w
    = \frac{1}{D}\,\mathbf{x}_v
      + \Bigl(1-\frac{1}{D}\Bigr)\overline{\mathbf{x}}_{2\text{-hop}} ,
    \label{eq:backtrack}
\end{equation}
because $v \in \mathcal{N}(u)$ for every $u \in \mathcal{N}(v)$: the walk can step
back. The ego node therefore receives weight exactly $1/D$, and it carries
\emph{undeflated} signal. A recursion that propagates only outwards describes the
\emph{non-backtracking} operator, not $L$ rounds of mean aggregation, and the
tree-like assumption does not license the substitution---the backtracking term is
present on an infinite tree. We therefore work with the true operator
$P = D^{-1}A$.

\begin{theorem}[Depth]
\label{thm:depth}
Let $q^{(L)}_j$ be the probability that an $L$-step simple random walk from $v$ on
the $D$-regular tree ends at distance $j$, and $N_j = D(D-1)^{j-1}$ ($N_0=1$) the
number of nodes at that distance. Then, in the noise-dominated regime,
\begin{equation}
    \frac{\mathrm{SNR}_L}{\mathrm{SNR}_{\mathrm{raw}}}
    \;=\;
    \frac{\bigl|\sum_j q^{(L)}_j \lambda^{j}\bigr|}
         {\sqrt{\sum_j \bigl(q^{(L)}_j\bigr)^2 / N_j}},
    \qquad \lambda = 1-(1+c)\rho .
    \label{eq:depth-exact}
\end{equation}
The label sequence along the walk is a two-state Markov chain with second
eigenvalue $\lambda$, which is why the depth-$j$ contribution is deflated by
$\lambda^{j}$ rather than by $\lambda^{L}$.
\end{theorem}

\begin{proof}
Deferred to Appendix~\ref{app:depth}.
\end{proof}

At $L=1$, $q_1 = 1$ and $N_1 = D$, so~\eqref{eq:depth-exact} reduces to
$|\lambda|\sqrt D$ and recovers Corollary~\ref{cor:critical}. For $L \geq 2$ it does
not factorise, and two consequences follow.

First, \textbf{the crossover is not depth-independent}. Solving
\eqref{eq:depth-exact} $= 1$ at $D=16$, $\pi_B = 0.3$ gives
$\rho^\star_1 = 0.525$, $\rho^\star_2 = 0.587$, $\rho^\star_3 = 0.553$,
$\rho^\star_4 = 0.576$, $\rho^\star_5 = 0.557$. Every value for $L\geq2$ exceeds the
single-layer threshold, so \textbf{depth buys tolerance to camouflage}.

The mechanism is that backtracking retains weight at distances $j < L$, which are
deflated by $\lambda^{j}$ rather than by $\lambda^{L}$; the ego term $j=0$
of~\eqref{eq:backtrack} is the extreme case, being undeflated entirely. It is worth
being precise about this, because the ego term is \emph{not} present at every layer:
the tree is bipartite, so $q^{(L)}_0 = 0$ for all odd $L$. At $L=3$ the ego weight is
zero, yet $\rho^\star_3 = 0.553$ still exceeds $\rho^\star_1$---there the gain comes
from mass returned to distance $1$ ($q^{(3)}_1 = 0.121$). The oscillation in $L$ is
this parity: even $L$ admits an ego term, odd $L$ does not.

Two qualifications matter. Deleting the ego term at $L=2$ and renormalising drops the
crossover from $0.587$ to $0.522$, so at even depths it does carry most of the effect.
And the benefit is specific to the camouflaged regime: backtracking \emph{costs}
separability when camouflage is absent, since the retained near-origin mass is
over-weighted noise. At $\rho=0$ the SNR ratios are $4.00, 11.49, 29.73, 72.91$ for
$L=1,\dots,4$ against $4.00, 15.49, 60.0, 232.4$ for a non-backtracking walk. Depth
under backtracking therefore trades peak separability at low $\rho$ for a higher
crossover---a bias--variance trade rather than a free gain, and one that favours
architectures able to tune the ego weight rather than inherit it at
$1/D$~\citep{zhu2020beyond,zhu2022how}.

Second, the harmful region is \emph{bounded above}. Only $|\lambda|$ enters, so at
$L=1$ aggregation is harmful precisely on the window
\begin{equation}
    \rho \;\in\;
    \left(\frac{1 - D^{-1/2}}{1+c},\;\; \frac{1 + D^{-1/2}}{1+c}\right)
    \cap (0,1] ,
    \label{eq:window}
\end{equation}
and helps again above it: a bot with $\rho \to 1$ attaches almost exclusively to
humans, so its neighbourhood is strongly \emph{anti}-correlated with its label and
informative once more, with the orientation reversed. We claim no novelty for this.
It is the second root of the same inequality
$\sqrt{D}\,|1-2\rho| > 1$ that Remark~\ref{rem:ma} attributes to
Ma et al.~\citep{ma2022homophily}, whose symmetry in $|p-q|$ those authors discuss
directly, and the resulting ``intermediate heterophily is worst'' phenomenon is the
mid-homophily pitfall named by Luan et al.~\citep{luan2023when}. We record it because
it identifies the dangerous adversarial strategy as \emph{intermediate} rather than
maximal camouflage, and because the upper edge exceeds $1$---so that the window
degenerates to a one-sided threshold and no ``helps again'' regime exists---whenever
$c < D^{-1/2}$. Among the relations measured in Section~\ref{sec:measuring}
($c = 0.369$), this holds for the sparser social ties (following, mention, quote,
reply, all $\tilde D \le 4$), but not for friendship or for the denser
content-coordination relations, where $D^{-1/2} < c$ and a finite upper edge
survives instead; Section~\ref{sec:r-estimate} reports that edge directly, under
the exact criterion, for shared URL and shared hashtag, and finds the measured
$\hat\rho$ still inside the harmful, non-reversed part of the window in both
cases.

\subsection{Implications for automated detection}
\label{sec:implications}

The three results above are statements about network structure, but they bear directly
on how automated accounts are detected in practice, and it is worth making the
connection explicit because the current literature addresses these mechanisms
architecturally rather than by measurement.

Detectors built on neighbourhood aggregation now routinely include a module that
down-weights or removes suspicious ties, a second aggregation channel for
dissimilar neighbours, or an explicit separation of an actor's own attributes from its
neighbours'~\citep{dou2020caregnn,he2025bothp,dang2026xhbot}. Each of these is
justified in the literature by an ablation: removing the module costs accuracy on a
benchmark. Our analysis supplies the missing account of \emph{why} they help and,
more usefully, \emph{when}.

Tie-discarding modules are the case analysed in Section~\ref{sec:achievable}, and the
analysis is less encouraging than the ablations suggest: discarding ties trades bias
against the number of ties left to average over, so its benefit is bounded by the quality of the underlying suspicion score, and on the networks we measure a score
based on attribute discrepancy alone is far too weak to restore inference
(Section~\ref{sec:measuring}). Separating an actor's own attributes from its
neighbours' is, by Theorem~\ref{thm:depth}, a way of making explicit and tunable a
term that iterated averaging already supplies at the fixed weight $1/D$---which
explains why the device helps most in exactly the regime, low degree and high
camouflage, where that fixed weight is least appropriate. And the finding that
camouflage differs by a factor approaching two across relation types within one
population implies that a detector pooling relations into a single adjacency is
averaging a usable signal with an unusable one; our thresholds apply per relation, and
suggest treating relation selection as a measurement question prior to model choice.

We draw one broader methodological conclusion. Because $\rho$, $c$, $D$ and the
attribute separation are all estimable from data a practitioner already has, the
question ``is relational information worth using here?'' can be answered before
fitting any model, and per relation rather than for a network as a whole. That is a
more informative starting point than an ablation, which reports only that a component
helped on one benchmark without indicating the conditions under which it would stop
helping.



\section{Numerical Validation}
\label{sec:validation}

Every closed form in Section~\ref{sec:theory} is checked against Monte-Carlo
simulation of Definition~\ref{def:ccsbm}. Our purpose here is verification of the
analysis, not a performance claim: the results below establish that the derivations
are correct, and say nothing on their own about accuracy on real bot graphs.
Section~\ref{sec:limitations} addresses that gap directly.

\paragraph{Protocol.}
Unless stated otherwise we use $\pi_B = 0.3$ ($c \approx 0.4286$),
$\|\Delta\| = 2$, $\sigma = 1$, and $D = 16$, drawing $4\times 10^5$ nodes per class
per configuration ($1.5 \times 10^6$ for the crossover sweep). Errors are computed
by exhaustive threshold search over the empirical distributions, so the reported
figures are attained minima rather than plug-in estimates.
Appendix~\ref{app:repro} records the remaining numerical choices.

\subsection{Deflation and error}

Table~\ref{tab:deflation} compares Proposition~\ref{prop:deflation} and
Theorem~\ref{thm:error} against simulation. Separation matches to within
$1.8\times 10^{-3}$ and the standard deviation---including the mixture term that a
naive calculation omits---to within $8\times 10^{-4}$ across the whole range.
Figure~\ref{fig:deflation} shows the same agreement graphically.

\begin{table}[t]
\centering
\caption{Theory versus simulation for signal deflation and balanced error
($\pi_B = 0.3$, $\|\Delta\|=2$, $\sigma=1$, $D=16$).}
\label{tab:deflation}
\scriptsize
\begin{tabular}{@{}ccccccc@{}}
\toprule
& & \multicolumn{2}{c}{separation $\delta$} & \multicolumn{2}{c}{$s_{\mathrm{B}}$} & error \\
\cmidrule(lr){3-4}\cmidrule(lr){5-6}
$\rho$ & $\eta$ & theory & sim. & theory & sim. & thy / sim. \\
\midrule
0.00 & 0.000 & 2.0000 & 2.0000 & 0.2500 & 0.2499 & 0.00003 / 0.00003 \\
0.10 & 0.043 & 1.7143 & 1.7142 & 0.2915 & 0.2918 & 0.00113 / 0.00160 \\
0.20 & 0.086 & 1.4286 & 1.4283 & 0.3202 & 0.3204 & 0.00925 / 0.01095 \\
0.30 & 0.129 & 1.1429 & 1.1441 & 0.3391 & 0.3391 & 0.03698 / 0.03895 \\
0.40 & 0.171 & 0.8571 & 0.8586 & 0.3500 & 0.3499 & 0.09785 / 0.09898 \\
0.50 & 0.214 & 0.5714 & 0.5696 & 0.3536 & 0.3528 & 0.19898 / 0.19974 \\
0.60 & 0.257 & 0.2857 & 0.2864 & 0.3500 & 0.3500 & 0.33734 / 0.33680 \\
\botrule
\end{tabular}
\end{table}

\begin{figure}[t]
    \centering
    \includegraphics[width=\linewidth]{figures/theory_deflation.png}
    \caption{Theory versus simulation. (a) Post-aggregation separation decays
    linearly in $\rho$ exactly as Proposition~\ref{prop:deflation} predicts.
    (b) Balanced error against the attribute-only baseline; the predicted critical ratio
    $\rho^\star = 0.464$ coincides with the crossing point.}
    \label{fig:deflation}
\end{figure}

\subsection{The critical ratio}

The central prediction is Corollary~\ref{cor:critical}. Solving
$\mathrm{SNR}(\rho) = \|\Delta\|/(2\sigma)$ numerically at $D=16$ gives
$\rho^\star = 0.4644$. Sweeping $\rho$ in steps of $0.01$ with $1.5\times 10^6$
samples per point, aggregation beats raw features at $\rho = 0.46$
($\varepsilon = 0.1545$ against $\varepsilon_{\mathrm{raw}} = 0.1587$) and loses at
$\rho = 0.47$ ($\varepsilon = 0.1648$). The measured crossover therefore lies in
$[0.46, 0.47]$, containing the predicted value.

Table~\ref{tab:critical} separates the two forms of the corollary. The exact
crossover differs from the closed form~\eqref{eq:critical} at moderate SNR, as
expected, since~\eqref{eq:critical} drops the mixture variance. Re-running the
exact solver in the noise-dominated regime it was derived for ($\|\Delta\| = 0.2$,
$\sigma = 1$) recovers the closed form to within $1.2\times10^{-3}$ at every degree,
confirming that the discrepancy is the stated approximation and not an error.

\begin{table}[t]
\centering
\caption{Critical camouflage ratio: closed form~\eqref{eq:critical} against the
exact SNR crossover, at moderate and low SNR.}
\label{tab:critical}
\small
\begin{tabular}{@{}cccc@{}}
\toprule
$D$ & closed form & exact ($r=2$) & exact (low SNR) \\
\midrule
4  & 0.3500 & 0.2604 & 0.3488 \\
9  & 0.4667 & 0.3913 & 0.4657 \\
16 & 0.5250 & 0.4644 & 0.5243 \\
25 & 0.5600 & 0.5102 & 0.5594 \\
64 & 0.6125 & 0.5807 & 0.6121 \\
\botrule
\end{tabular}
\end{table}

Figure~\ref{fig:critical} plots both curves against degree. The tolerated
camouflage rises steeply for sparsely connected accounts and flattens as $D$ grows,
which has a concrete operational reading: the accounts most vulnerable to
camouflage are those with few connections, and a deployment that aggregates over
small neighbourhoods is fragile in exactly the regime where bots are cheapest to
create.

\begin{figure}[t]
    \centering
    \includegraphics[width=0.72\linewidth]{figures/theory_critical_ratio.png}
    \caption{Critical camouflage ratio against neighbourhood size. Below the lower
    curve, averaging over alters improves on attributes alone; between the curves it is
    actively harmful; above the upper edge $(1+D^{-1/2})/(1+c)$ the neighbourhood is
    strongly anti-correlated with the label and informative again with the orientation
    reversed. The closed form is the noise-dominated limit of the exact crossover, and
    the upper edge exceeds $1$---so that no reversal regime exists---whenever
    $c < D^{-1/2}$.}
    \label{fig:critical}
\end{figure}

\subsection{The tie-scoring rule}

Lemma~\ref{lem:energy} predicts expected squared discrepancies of $2d\sigma^2$ and
$\|\Delta\|^2 + 2d\sigma^2$. At $d = 8$, $\sigma = 1$, $\|\Delta\| = 2$ these are
$16$ and $20$; simulation over $2\times10^5$ edge samples gives $16.005$ and
$20.005$, an error below $0.03\%$.

For Theorem~\ref{thm:sgtr} we verify the closed
form~\eqref{eq:rhostarsgtr} by round-trip: substituting
$\rho^\star_{\mathrm{prune}}$ back into $\rho_{\mathrm{eff}}$ must return
$\rho^\star$. The maximum absolute deviation over
$\bar e_c \in \{0.25, 0.5, 0.75, 0.9\}$ is $1.1\times 10^{-16}$, i.e.\ machine
precision. Table~\ref{tab:sgtr} reports the resulting tolerance gains, and
Figure~\ref{fig:sgtr} plots the full curve.

\begin{table}[t]
\centering
\caption{Effect of the spectral rule attenuation on tolerated camouflage under the
composition-only form~\eqref{eq:rhostarsgtr} ($D=16$, $\pi_B=0.3$, baseline
$\rho^\star = 0.525$). These values ignore both the effective-degree cost and the
correct $\eta_{\mathrm{eff}}$; Section~\ref{sec:validation-ext} gives the figures
we stand behind. The table is retained only to show the size of the two corrections.}
\label{tab:sgtr}
\small
\begin{tabular}{@{}ccc@{}}
\toprule
$\bar e_c$ & $\rho_{\mathrm{eff}}$ at $\rho = 0.4$ & $\rho^\star_{\mathrm{prune}}$ \\
\midrule
0.00 & 0.4000 & 0.5250 \\
0.25 & 0.3333 & 0.5957 \\
0.50 & 0.2500 & 0.6885 \\
0.75 & 0.1429 & 0.8155 \\
0.90 & 0.0625 & 0.9170 \\
\botrule
\end{tabular}
\end{table}

\begin{figure}[t]
    \centering
    \includegraphics[width=0.78\linewidth]{figures/theory_pruning_tolerance.png}
    \caption{Tolerated camouflage against the attenuation applied to camouflage ties.
    The solid curve is the corrected analysis (Proposition~\ref{prop:deff}, with
    $\eta_{\mathrm{eff}}$ as in~\eqref{eq:pruned-condition}); the dotted curve is the
    composition-only form~\eqref{eq:rhostarsgtr}. The correction is \emph{upward} over
    the plotted range: treating the majority side correctly matters more than the
    effective-degree penalty. The curve is truncated at $\bar e_c = 0.9$ deliberately.
    Equation~\eqref{eq:deff} replaces the random retained count by its mean, which is
    defensible while many ties survive but is exactly the substitution
    Remark~\ref{rem:no-ceiling} shows to be invalid as $\bar e_c \to 1$; that limit is
    governed instead by Proposition~\ref{cor:ceiling}, which says perfect pruning never
    loses. The dash-dotted line marks where half of the \emph{strategic actors} still
    retain enough ties to benefit; majority actors are barely affected.}
    \label{fig:sgtr}
\end{figure}

\subsection{The extensions: achievable pruning and depth}
\label{sec:validation-ext}

The effective-degree formula~\eqref{eq:deff} matches Monte-Carlo variance to
within $0.5\%$ for $\theta \geq 0.5$. At $\theta = 0$ it is off by $7.3\%$
($8.00$ predicted against $7.41$ measured at $\rho=0.5$, $D=16$), and the cause is
not sampling: it is Jensen's inequality on $\mathbb{E}[1/K]$ for the random retained
count $K$, exactly the effect quantified in Remark~\ref{rem:no-ceiling}. Under
perfect pruning the variance-equivalent degree is $1.29$ rather than the $1.00$ that
$\mathbb{E}[K]$ predicts at $\rho = 1-1/D$, $D=16$, and $35.6\%$ of nodes retain no
neighbour at all; the proportion still benefiting falls to one half at
$\rho = 0.897$ ($D=16$) and $0.974$ ($D=64$).

Correcting $\eta_{\mathrm{eff}}$ as in~\eqref{eq:pruned-condition} raises the
tolerances relative to the $c\,\rho_{\mathrm{eff}}$ substitution: at $D=16$,
$\pi_B=0.3$ we obtain $0.609$, $0.715$, $0.847$ and $0.928$ for
$\bar e_c = 0.25, 0.5, 0.75, 0.9$, against $0.594$, $0.679$, $0.788$ and $0.869$
under the incorrect form. The error was therefore pessimistic, and by up to six
percentage points.

For depth we evaluate~\eqref{eq:depth-exact} directly, iterating the
distance recursion for the walk profile and summing the weighted contributions. The weight profile $\sum_j (q^{(L)}_j)^2/N_j$ exceeds the
non-backtracking value $D^{-L}$ by factors of $1.94$, $4.63$, $12.3$ and $35.0$ at
$L = 2,\dots,5$, which is the quantitative size of the term
that~\eqref{eq:backtrack} restores. The resulting crossovers are
$\rho^\star_1 = 0.5250$, $\rho^\star_2 = 0.5868$, $\rho^\star_3 = 0.5528$,
$\rho^\star_4 = 0.5756$ and $\rho^\star_5 = 0.5567$: every multi-layer value lies
above the single-layer one, so depth increases the tolerated camouflage, and the
crossover is \emph{not} depth-independent.

\subsection{Measuring the camouflage ratio on real graphs}
\label{sec:measuring}

The theory is parameterised by $\rho$, which has so far been a modelling device. We
now measure it. TwiBot-20 and TwiBot-22 require author approval for access, which we
did not have. We therefore use \textbf{MGTAB}~\citep{shi2023mgtab}, an openly
distributed Twitter dataset of $10{,}199$ expert-annotated accounts
($2{,}748$ automated, $7{,}451$ genuine; $c = 0.369$) in which \emph{every} account
carries a label, so that all ties are between labelled actors and the measurement is
not diluted by unlabelled alters. MGTAB records seven relations on the same
population, which lets us separate the social layer (following, friendship, mention,
reply, quote) from a content-coordination layer (shared URLs, shared hashtags). We also use the
two review networks of Dou et al.~\citep{dou2020caregnn} as a contrast case,
reported in the prose below rather than in Table~\ref{tab:real}; these are not social
ties and we draw no structural conclusion from them, but they carry the attribute
information needed to evaluate the tie-scoring bound of
Section~\ref{sec:achievable}, which MGTAB's inflated attribute separation does not
support.

\paragraph{A negative result on the canonical benchmark.}
It is worth recording why we do not use cresci-2015 or cresci-2017, the
longest-standing bot benchmarks. In the openly distributed preprocessed release,
cresci-2017 contains \emph{no} interpersonal ties at all: all $6{,}637{,}615$ of its
edges are account-to-tweet postings. cresci-2015 does contain following and
friendship ties, but only $0.41\%$ of the $1.29$ million accounts they span are
labelled, and the labelled-to-labelled subgraph reduces to $4{,}275$ ties with the
mixing pattern $3{,}823$ genuine--genuine, $373$ automated--automated and just
$\mathbf{79}$ cross-class. Only $698$ of $3{,}351$ automated accounts ($20.8\%$) have
any labelled tie, the median such account has degree $0$, and the median camouflage
ratio is therefore $0$. A quantity resting on $79$ ties cannot support a
distributional claim, and we obtained the same counts from a second, independently
preprocessed release of the same data, so this is a property of the dataset rather
than of one pipeline. We report it because it explains a gap in the literature: the
benchmark most used for bot detection cannot be used to measure the structural
mechanism the field attributes to bots.

For each relation we compute $\rho$ as the mean over automated accounts of the
fraction of their alters that are genuine, and $\eta$ symmetrically.
Table~\ref{tab:real} reports the result.

\begin{sidewaystable}
\centering
\small
\caption{Camouflage measured on MGTAB ($10{,}199$ Twitter accounts, all labelled;
$c = 0.369$). $\hat\rho$ is the mean genuine share of an automated account's alters,
$\tilde\rho$ its median, $\tilde D$ the median degree over automated accounts having
at least one tie of that relation. The last two columns evaluate the \emph{exact}
criterion of Theorem~\ref{thm:error} at $r = 2.32$ (Section~\ref{sec:r-estimate}),
not the noise-dominated window~\eqref{eq:window}: $\rho_{\mathrm{cross}}$ is the
crossover and the final column is
$\mathrm{SNR}_{\mathrm{agg}}/\mathrm{SNR}_{\mathrm{raw}}$ at the measured $\hat\rho$,
so values below $1$ mean the relation is worse than using attributes alone. The
penultimate column evaluates that ratio with the model-implied $\eta = c\rho$; the final
column substitutes the \emph{measured} $\hat\eta$, which on this population is far
smaller because automated accounts have much lower degree
(Section~\ref{sec:balance-check}). The verdict is unchanged either way; the magnitude is
not.}
\label{tab:real}
\begin{tabular}{lrrrrrrrr}
\toprule
Relation & ties & $\hat\rho$ & $\tilde\rho$ & $\hat\eta$ & $\tilde D$
 & $\rho_{\mathrm{cross}}$
 & SNR$_{\mathrm{agg}}$/SNR$_{\mathrm{raw}}$
 & \,''\,($\hat\eta$) \\
\midrule
\multicolumn{9}{@{}l}{\emph{social layer}}\\
following   & 231{,}829 & 0.969 & 1.00 & 0.041 & 4  & 0.252 & 0.51 & 0.02 \\
friendship  & 327{,}035 & 0.989 & 1.00 & 0.033 & 8  & 0.372 & 0.79 & 0.06 \\
mention     & 107{,}794 & 0.994 & 1.00 & 0.022 & 4  & 0.252 & 0.57 & 0.03 \\
quote       & 20{,}536  & 0.999 & 1.00 & 0.014 & 2  & 0.121 & 0.42 & 0.02 \\
reply       & 20{,}700  & 0.991 & 1.00 & 0.009 & 1  & ---   & 0.28 & 0.00 \\
\;\;union   & 331{,}251 & 0.989 & 1.00 & 0.032 & 8  & 0.372 & 0.79 & 0.06 \\
\midrule
\multicolumn{9}{@{}l}{\emph{content-coordination layer}}\\
shared URL     & 263{,}800 & 0.595 & 0.62 & 0.195 & 12 & 0.432 & 0.44 & 0.51 \\
shared hashtag & 300{,}000 & 0.627 & 0.65 & 0.279 & 53 & 0.585 & 0.71 & 0.47 \\
\botrule
\end{tabular}
\end{sidewaystable}

\subsection{Estimating the attribute separation}
\label{sec:r-estimate}

The criterion depends on $r = \|\Delta\|/\sigma$, and this quantity is easy to
overstate. The obvious plug-in---$\|\Delta\|$ from the class-mean difference and
$\sigma$ from the mean per-dimension variance---gives $r = 13.0$ on MGTAB's $788$
attributes. That figure is not usable: it implies an attribute-only error of
$\Phi(-6.5) \approx 4\times10^{-11}$, i.e.\ that the attributes classify the
population essentially perfectly and the question of whether the network helps does
not arise. A held-out linear discriminant on the same attributes attains a balanced
error of $0.123$. The discrepancy is the isotropy assumption: the per-dimension
standard deviations span a factor of $28$, so a single pooled $\sigma$ is not
meaningful, while $\|\Delta\|^2$ accumulates small differences over many correlated
coordinates.

We therefore calibrate $r$ from the achievable error rather than from the moments,
taking $r = -2\Phi^{-1}(0.123) = 2.32$, and we use that value throughout
Table~\ref{tab:real}. Resampling the split moves $r$ materially---over twelve random splits we obtain $r = 2.24 \pm 0.07$, a range of $0.22$---so we report the mean rather than a single split. The verdicts are insensitive to it: re-solving the exact criterion across that whole range leaves every relation harmful, the largest SNR ratio moving only between $0.79$ and $0.82$. This matters: at $r = 13$ the mixture term
$\rho(1-\rho)\|\Delta\|^2$ reaches $1.67\sigma^2$ and the noise-dominated
form~\eqref{eq:window} is invalid, whereas at $r = 2.32$ it is $0.05\sigma^2$ and the
approximation is defensible. Any application of these thresholds should report $r$
and how it was obtained.

\paragraph{On this population, every relation is one the network harms.}
The last column of Table~\ref{tab:real} is below $1$ for all seven recorded relations, and for their social union: under
the exact criterion, neighbour averaging on MGTAB is worse than using attributes alone
on \emph{every} one of the seven recorded relations. The size of the shortfall depends
on which $\eta$ is used: with the model-implied $\eta = c\rho$ the factors run from
$1.3$ to $3.6$, whereas substituting the measured $\hat\eta$---which is the honest
choice on a population whose classes differ in degree by a factor of five---they run from
$2$ on the content relations to more than an order of magnitude on the social ones. We
report both, because the gap between them is itself the degree effect of
Section~\ref{sec:balance-check} and not a modelling detail to be hidden. We report this
as the principal empirical finding, and we note that an earlier version of this
analysis reached the opposite conclusion for the social relations by applying the
noise-dominated window~\eqref{eq:window} at $r = 13$; the ``informative again with
reversed sign'' regime that produced is an artefact of using an $r \to 0$
approximation at large $r$, and it does not survive the exact criterion. Where the
harmful interval does have an upper edge at $r = 2.32$---shared URL $(0.432, 0.995)$
and shared hashtag $(0.585, 0.868)$---the measured $\hat\rho$ falls inside it.

\paragraph{The layers still differ, in degree of harm and in mechanism.}
What survives, and is not an artefact, is the contrast between the two layers.
On every social relation the automated accounts are almost perfectly camouflaged:
$\tilde\rho = 1.00$, so the median automated account has \emph{no} automated alter,
and automated--automated ties number $212$ in the following-and-friendship union
against $40{,}263$ automated--genuine ties. On the content relations the same accounts
behave in the opposite way: $\hat\rho$ falls to $0.60$--$0.63$ and
automated--automated ties rise to $15{,}963$ (shared URL) and $28{,}203$ (shared
hashtag). These accounts avoid one another where ties are visible and converge where
content is shared. The ratio $\hat\rho_{\text{social}}/\hat\rho_{\text{content}}$ is
$1.68$, and the resulting harm differs correspondingly: the SNR ratio ranges from
$0.79$ on friendship to $0.44$ on shared URLs.

\paragraph{What attribute-based tie scoring achieves.}
Section~\ref{sec:achievable} bounds the tolerance attainable by a threshold rule on
attribute discrepancy; we report the values here, calibrating $r$ on these networks the
same way as on MGTAB, from the balanced error of a held-out linear discriminant. That
matters, because the moment-based plug-in disagrees with it in \emph{both} directions:
on YelpChi it gives $1.17$ against a held-out $1.05$, and on Amazon $2.04$ against
$2.65$. Reporting a single method for all three datasets is the only defensible choice.

On YelpChi \texttt{homo} ($c=0.170$, $\tilde D=158$, $d=32$) the optimum moves the
tolerance from $0.7867$ to $0.7870$; on Amazon \texttt{homo} ($c=0.074$, $\tilde D=145$,
$d=25$) from $0.8539$ to $0.8808$. Both remain below the measured $\hat\rho$ of $0.805$
and $0.897$, so attribute-discrepancy scoring does not restore the inference on either
network---but we note the margins are not large, $0.018$ and $0.016$ respectively, and
that the Amazon margin halves under the corrected calibration. A modestly better score
would close it. The optimum is interior: it falls to zero as $\bar e_w \to 1$, because a
threshold rule that discards nearly every tie leaves nothing to average, and it is
stable under enlarging the search domain, which the scale-invariant soft-weight model
of~\eqref{eq:deff} is not.

These two tolerances are computed in the noise-dominated regime, unlike the MGTAB
figures in Table~\ref{tab:real}, which use the exact criterion. At these networks'
$\hat\rho$ the mixture term is $0.22\sigma^2$ (YelpChi) and $0.65\sigma^2$ (Amazon), so
the approximation is loose; the exact criterion gives slightly lower tolerances
($0.780$ and $0.840$) and therefore strengthens rather than weakens the conclusion. We
report the noise-dominated values because they are the ones the closed form yields, and
flag the discrepancy rather than let it pass silently.

\paragraph{Multiplexity, not density, decides usability.} These relations are recorded
on the \emph{same} $10{,}199$ accounts, and both $\hat\rho$ and the crossover differ
substantially between them. A procedure that pools relations into a single adjacency
is therefore averaging quantities with different thresholds, which is the concrete
form of the concern raised by V\"or\"os and Snijders~\citep{voros2017cluster} and by
Snijders et al.~\citep{snijders2013model}: the quantity is a property of a relation,
not of a population.

\paragraph{The degree-corrected identity is exact, once the estimators match.}
\label{sec:balance-check}
On the social union, the identity in its degree-corrected form
$\eta = c\rho\,(\bar d_B/\bar d_H)$ predicts $0.0647508$ against a measured
$0.0647508$: agreement to machine precision. This is not a coincidence but a
consequence of what~\eqref{eq:balance} is---an accounting identity over edge
endpoints---and it holds on \emph{any} graph provided $\rho$, $\eta$ and the mean
degrees are all estimated as ratios of sums, matching the sums the derivation counts.
The uncorrected form $\eta = c\rho$ predicts $0.365$, so the degree factor
$\bar d_B/\bar d_H = 0.177$ is doing essential work here: automated accounts have mean
degree $14.8$ against $83.5$ for genuine ones.

We flag the estimator sensitivity because it caught us out. Computing $\hat\eta$ as a
\emph{mean of per-node ratios}, while predicting it from a \emph{ratio of means},
yields $0.032$ against a predicted $0.067$ and the appearance of a factor-of-two
failure. The discrepancy is entirely an artefact of mixing the two estimators, com\-pounded
by averaging degrees over connected nodes only. Since a mean of ratios weights
low-degree actors equally with hubs, the two estimators diverge exactly in proportion
to within-class degree skew, which on this population is severe. The practical
guidance is that~\eqref{eq:balance} is a constraint an analyst can \emph{check} on
their own data, and a failed check indicates an estimator mismatch or a violated
regularity assumption rather than a violated identity.

\paragraph{Caveat.}
The window is derived for a CSBM with regular degrees, isotropic Gaussian features,
and a single relation. Real graphs violate all three: degrees are heavy-tailed,
features are neither Gaussian nor isotropic, and the relations interact. The
numbers above should be read as locating these graphs approximately, not as a
certified verdict, and the agreement of $\eta$ with $c\rho$ on YelpChi is the one
quantity here that constitutes a genuine test of the theory rather than an
application of it. Measuring $\rho$ on TwiBot remains the natural next step and
requires only data access. The estimator itself is a direct count, and on synthetic
graphs with a planted camouflage ratio it recovers that ratio exactly and reproduces
$\eta = c\rho$ to machine precision under regular degrees, so the measurement is
limited by data availability rather than by method.



\section{Discussion}
\label{sec:discussion}

\subsection{From heuristic to criterion}

The practical value of Corollary~\ref{cor:critical} is that it converts a design
intuition into a quantity an operator can estimate. The inputs are the mean
neighbourhood size $D$, the bot prevalence $\pi_B$, and the feature signal-to-noise
ratio $\|\Delta\|/\sigma$---all measurable on a deployed system---and the output is
a threshold on camouflage beyond which the graph is a liability. This reframes a
question usually settled by ablation (``does the graph help on this benchmark?'')
as one that can be answered per-cohort: for hub accounts with large $D$ the graph
remains valuable at camouflage levels that make it actively harmful for sparse
accounts. A detector that applies one aggregation policy to both is leaving
accuracy on the table at one end and incurring damage at the other.

The analysis also isolates an under-appreciated coupling. Because
$\eta = c\rho$, bot prevalence enters the critical ratio directly: a platform whose
bot population doubles does not merely face more bots, it faces a detector whose
tolerance to each bot's camouflage has fallen. Class imbalance is normally treated
as a training-time nuisance to be fixed by resampling; \eqref{eq:critical} shows it
is also a structural robustness parameter, and resampling does not touch it.

\subsection{What the theory does and does not license}

We are deliberate about the boundaries of these claims.

The results characterise \emph{one round of mean aggregation} in a stylised
generative model. They explain why homophilic message passing fails under
camouflage and why the three mitigations help, and they give the correct
qualitative orderings---more pruning is better, larger neighbourhoods are safer,
higher prevalence is more dangerous. They do not predict the accuracy of a trained
multi-layer network on a real bot graph, and no number in
Section~\ref{sec:validation} should be read as such. Our simulations verify
algebra; they are not empirical evidence about Twitter or any other platform.

Three specific gaps deserve naming. First, the single-relation restriction is real: our measurements show camouflage
differs by a factor approaching two across relations of the same graph
(Table~\ref{tab:real}), and a multi-relational treatment would need to model how a
detector pools them. Second, although Section~\ref{sec:depth} extends the analysis
to $L$ layers, it does so in the noise-dominated regime and on a regular tree; a
treatment that keeps the mixture terms at depth, and that handles the finite graphs
where receptive fields saturate, is open. Third, Section~\ref{sec:achievable} bounds
the achievable pruning quality from the feature distributions but treats the scorer
as a fixed threshold rule; bounding what a \emph{learned} scorer attains, as a
function of sample size and hypothesis class, is the natural next step and is what
our real-data measurements suggest matters most and would
turn Theorem~\ref{thm:sgtr} from a conditional into an unconditional guarantee.\subsection{Limitations of the generative model}
\label{sec:limitations}

The camouflage-CSBM makes four simplifications, each of which trades realism for
tractability.

\emph{Isotropic, equal-covariance features.} Real account features are neither
isotropic nor identically dispersed across classes. Anisotropy would replace
$\|\Delta\|/\sigma$ with a Mahalanobis distance throughout; the structure of the
results survives, but $\rho^\star$ becomes direction-dependent.

\emph{Fixed neighbourhood size.} Social graphs are heavy-tailed. Since $\rho^\star$
is concave in $D$, a mixture of degrees is \emph{not} well summarised by its mean:
applying~\eqref{eq:critical} at the average degree overstates robustness for the
low-degree majority. Per-degree stratification, which the theory directly supports,
is the appropriate remedy.

\emph{Independent neighbour sampling.} We ignore triadic closure and community
structure, both prominent in real follow graphs. Correlated neighbourhoods reduce
the effective sample size in the averaging step, which inflates the variance
in~\eqref{eq:variances} and lowers $\rho^\star$; our thresholds are therefore
optimistic.

\emph{A single homogeneous relation.} Production detectors operate on
heterogeneous graphs with several edge types. The analysis extends per-relation,
and the interesting question---how an adversary should allocate a fixed camouflage
budget across relations---follows naturally from~\eqref{eq:critical} but is beyond
the present scope.

\subsection{An adversarial reading}

Finally, the theory can be read from the attacker's side, which we think is the
more useful test of it. Corollary~\ref{cor:critical} tells an adversary exactly how
much camouflage is needed to neutralise a homophilic detector, and
Theorem~\ref{thm:sgtr} tells them how much more is needed against a pruning
defence. That the achievable tie score is far weaker than a flawless one on the networks we
measure---moving the tolerance by less than $0.01$ on one of them---means no realistic
pruning rule is a permanent fix; the defence buys operating range, and an adversary
willing to spend ties can exceed it. Treating $\rho$ as a
strategic variable rather than a nuisance parameter---and analysing the resulting
game between detector and adversary---strikes us as the most promising direction
this line of work opens.

\section{Conclusion}

We close by restating the priority position plainly, because it determines how this
paper should be cited. That a crossover exists---a camouflage level beyond which mean
aggregation is worse than discarding the graph---is due to Ma
et al.~\citep{ma2022homophily}, and the fact that the harmful region is a bounded
window rather than a half-line is the second root of their condition, discussed under
the name mid-homophily pitfall by Luan et al.~\citep{luan2023when}. What we add is the
class-imbalanced adversarial case: the edge-balance identity $\eta = c\rho$, which
makes bot prevalence a robustness parameter and which we confirm on real data; a
mixture-variance correction worth $11.5\%$ of the threshold; the observation that
pruning trades bias against effective degree, so that a bound on achievable edge-score
quality is needed before any tolerance claim means anything; and the finding that mean
aggregation's backtracking term acts as an implicit residual connection, so that the
crossover rises with depth rather than being invariant to it.

\label{sec:conclusion}

We have given an analytical account of relation camouflage in GNN-based bot
detection. Modelling camouflage as an adversarial parameter in a contextual
stochastic block model, we proved that mean aggregation deflates the class signal
by exactly $1-(1+c)\rho$, derived the resulting Bayes error, and obtained a
critical camouflage ratio $\rho^\star = (1-D^{-1/2})/(1+c)$ beyond which a
homophilic GNN is provably worse than ignoring the graph---a result that recovers
the known $1/\sqrt{D}$ separability gain of graph convolution at $\rho = 0$. All closed forms were verified numerically, with
the predicted critical ratio falling inside the measured crossover bracket.

The contribution is not a new detector but a criterion. Practitioners currently
decide whether to trust a graph-based detector by running an ablation on a
benchmark; the results here let them instead estimate three measurable quantities
and compute the answer for their own deployment, per account cohort. Extending the
analysis to multiple layers, to heterogeneous relations, and to a learned rather
than assumed pruning quality would make that criterion sharper still.



\backmatter

\section*{Statements and Declarations}

\subsection*{Funding}
The authors declare that no funds, grants, or other support were received during the
preparation of this manuscript.

\subsection*{Conflict of interest/Competing interests}
The authors declare that they have no known competing financial interests or personal
relationships that could have appeared to influence the work reported in this paper.

\subsection*{Ethics approval and consent to participate}
Not applicable. This study analyses publicly released, third-party, de-identified
social-media account data; it did not involve new data collection from human
participants.

\subsection*{Consent for publication}
Not applicable.

\subsection*{Data availability}
All data analysed in this study are third-party datasets that are publicly available
and were not generated by the authors. The Twitter account data are from MGTAB
\citep{shi2023mgtab}, distributed by its authors under a Creative Commons
Attribution--NonCommercial--NoDerivatives 4.0 licence; we analysed the released
tensors without modification and redistribute no derivative of them. The two review
networks are those of \citet{dou2020caregnn} and are openly downloadable. The
cresci-2015 and cresci-2017 collections referred to in
Section~\ref{sec:measuring} are likewise public. TwiBot-20 and TwiBot-22 require
permission from their authors, which we did not obtain and which is why they are not
analysed here.

\subsection*{Materials availability}
Not applicable.

\subsection*{Code availability}
Code to reproduce the reported simulations, closed-form checks, and MGTAB
measurements is available at \url{https://github.com/vinhqdang/XHBot}
(directory \texttt{paper2/theory}).

\subsection*{Author contribution}
Quang-Vinh Dang: conceptualization, methodology, formal analysis, software,
validation, writing -- original draft, writing -- review and editing. Phuong-Lan
Nguyen: methodology, formal analysis, investigation, validation, writing -- review
and editing.

\begin{appendices}

\section{Deferred proofs}
\label{app:proofs}

This appendix gives in full the proofs that are compressed or omitted in the main
text. Throughout, $\pi_B$ is the bot prevalence, $c = \pi_B/(1-\pi_B)$,
$\eta = c\rho$, $\lambda = 1-(1+c)\rho$, features are $\mathbf{x}_v \sim
\mathcal{N}(\boldsymbol\mu_{y_v}, \sigma^2 I_d)$ with
$\Delta = \boldsymbol\mu_B - \boldsymbol\mu_H$, and $D$ is the neighbourhood size.

\subsection{The edge-balance identity}
\label{app:balance}

\begin{lemma}
Let camouflage edges be those joining a bot to a human, and suppose each bot directs
a fraction $\rho$ of its $D$ edges to humans. Then the expected fraction of a
human's neighbourhood that is composed of bots is $\eta = c\rho$.
\end{lemma}

\begin{proof}
Let $n$ be the number of nodes, so there are $\pi_B n$ bots and $(1-\pi_B)n$ humans.
Each bot contributes $\rho D$ camouflage edge-endpoints, giving
$\pi_B n \rho D$ endpoints in total. Every such endpoint lands on a human, and the
humans collectively expose $(1-\pi_B) n D$ endpoints. Hence the expected bot fraction
of a human neighbourhood is
\[
    \eta \;=\; \frac{\pi_B n \rho D}{(1-\pi_B) n D}
        \;=\; \frac{\pi_B}{1-\pi_B}\,\rho \;=\; c\rho .
\]
The identity is a conservation statement about edge endpoints and is independent of
how bots choose \emph{which} humans to attach to; only the aggregate count matters.
It presumes $\eta \leq 1$, i.e.\ $\rho \leq 1/c$, which is automatic when
$\pi_B \leq 1/2$.
\end{proof}

The asymmetry is the point: the adversary sets $\rho$ directly but induces $\eta$ on
the other class, and the two coincide only when the classes are balanced ($c=1$).

\subsection{Proof of the depth result (Theorem~\ref{thm:depth})}
\label{app:depth}

\begin{proof}
Project onto $\hat{\mathbf{u}} = \Delta/\|\Delta\|$. Write $P = D^{-1}A$ for the
simple-random-walk operator on the $D$-regular tree, so that $L$ rounds of mean
aggregation produce $\mathbf{h}^{(L)} = P^{L}\mathbf{x}$. Crucially $P$ is
\emph{not} supported on the $L$-sphere: since $v \in \mathcal{N}(u)$ whenever
$u \in \mathcal{N}(v)$, the walk backtracks, and $P^L e_v$ places mass at every
distance $j \le L$ of the same parity as $L$, together with $j=0$ when $L$ is even.

\emph{Weights.} By symmetry all nodes at distance $j$ carry equal weight, so if
$q^{(L)}_j$ denotes the probability that an $L$-step walk ends at distance $j$, each
such node has weight $q^{(L)}_j/N_j$ with $N_j = D(D-1)^{j-1}$ and $N_0 = 1$. The
distance process is the birth--death chain that steps outward with probability
$(D-1)/D$ and inward with probability $1/D$ from any $j\ge1$, and outward with
certainty from $j=0$; iterating it gives $q^{(L)}$.

\emph{Signal.} Along a path in the tree the labels form a two-state Markov chain with
transition matrix $\begin{pmatrix} 1-\rho & \rho\\ \eta & 1-\eta\end{pmatrix}$, whose
eigenvalues are $1$ and $\lambda = 1-\rho-\eta = 1-(1+c)\rho$. Hence for a node at
distance $j$,
$\mathbb{E}[s_w \mid y_v = B] - \mathbb{E}[s_w \mid y_v = H] = \lambda^{j}\|\Delta\|$,
and summing against the weights gives
$\Delta_L = \|\Delta\|\sum_j q^{(L)}_j \lambda^{j}$. Note this is \emph{not}
$\lambda^L\|\Delta\|$: the $j<L$ terms are deflated less, and the $j=0$ term not at
all.

\emph{Noise.} For independent $\mathcal{N}(0,\sigma^2)$ noise,
\begin{align*}
  \operatorname{Var}\bigl(\mathbf{h}^{(L)}_v\bigr)
  &= \sigma^2\bigl\|P^Le_v\bigr\|^2
  = \sigma^2\sum_j N_j \bigl(q^{(L)}_j/N_j\bigr)^2 \\
  &= \sigma^2\sum_j \bigl(q^{(L)}_j\bigr)^2 / N_j .
\end{align*}
Dividing signal by noise and normalising by
$\mathrm{SNR}_{\mathrm{raw}} = \|\Delta\|/(2\sigma)$ gives~\eqref{eq:depth-exact};
the absolute value is required because $\sum_j q_j\lambda^j$ can be negative.
At $L=1$, $q_1=1$ and $N_1=D$, recovering $|\lambda|\sqrt{D}$.
\end{proof}

\begin{remark}[Why the naive recursion fails, and by how much]
Propagating the deflation as $\Delta_{l+1} = \lambda\Delta_l$ and the variance as
$V_{l+1} = V_l/D$ describes the \emph{non-backtracking} walk, and yields
$\mathrm{SNR}_L/\mathrm{SNR}_{\mathrm{raw}} = (\lambda\sqrt D)^L$ and a
depth-independent crossover. Both conclusions are artefacts. The true weight
concentration $\sum_j (q^{(L)}_j)^2/N_j$ exceeds $D^{-L}$ by factors of $1.94$,
$4.63$, $12.3$ and $35.0$ at $L = 2,3,4,5$ for $D = 16$, and the crossover moves from
$0.525$ at $L=1$ to $0.587$ at $L=2$. The direction is intuitive once stated: mass retained at distance $j<L$ is deflated
only by $\lambda^{j}$, so backtracking raises the crossover. It does so at a cost,
however: at $\rho=0$ the same retained mass is over-weighted noise, and the
non-backtracking walk attains a strictly higher SNR ($60.0$ against $29.7$ at $L=3$,
$D=16$). Backtracking trades separability in the homophilous regime for tolerance in
the camouflaged one. A
simulation that averages over the $D^L$ leaves of a tree reproduces the naive
recursion exactly and therefore cannot detect this error; one must simulate $P^L$.
\end{remark}

\subsection{Exact law of the edge score}
\label{app:score}

\begin{lemma}
For an edge $(u,v)$ let $s_{uv} = \|\mathbf{x}_u - \mathbf{x}_v\|^2$. Then
$s \sim 2\sigma^2\chi^2_d$ if $y_u = y_v$, and
$s \sim 2\sigma^2\chi^2_d(r^2/2)$ with $r = \|\Delta\|/\sigma$ otherwise. In
particular $\mathbb{E}[s \mid \text{camouflage}] - \mathbb{E}[s \mid \text{within}]
= \|\Delta\|^2$, and the ratio of the gap to the within-class standard deviation is
$\kappa = r^2/(2\sqrt{2d})$.
\end{lemma}

\begin{proof}
Write $\mathbf{z} = \mathbf{x}_u - \mathbf{x}_v$. If $y_u = y_v$ then
$\mathbf{z} \sim \mathcal{N}(0, 2\sigma^2 I_d)$, so
$\|\mathbf{z}\|^2/(2\sigma^2) \sim \chi^2_d$. If $y_u \neq y_v$ then
$\mathbf{z} \sim \mathcal{N}(\pm\Delta, 2\sigma^2 I_d)$ and
$\|\mathbf{z}\|^2/(2\sigma^2)$ is noncentral chi-square with $d$ degrees of freedom
and noncentrality $\|\Delta\|^2/(2\sigma^2) = r^2/2$.

For the moments, $\mathbb{E}[\chi^2_d] = d$ and
$\mathbb{E}[\chi^2_d(\nu)] = d + \nu$, so the means are $2\sigma^2 d$ and
$2\sigma^2 d + \|\Delta\|^2$, giving the stated gap. Since
$\operatorname{Var}(\chi^2_d) = 2d$, the within-class variance of $s$ is
$4\sigma^4\cdot 2d$, whose square root is $2\sigma^2\sqrt{2d}$; dividing the gap by
it gives $\kappa = \|\Delta\|^2/(2\sigma^2\sqrt{2d}) = r^2/(2\sqrt{2d})$.

The $d^{-1/2}$ dependence is the concentration phenomenon: the gap between the two
score distributions grows like $\|\Delta\|^2$ while their spread grows like
$\sqrt{d}$, so a fixed feature separation becomes progressively harder to detect
from raw squared distances as the ambient dimension increases.
\end{proof}

\section{Computational details}
\label{app:repro}

Every closed form reported in the paper is checked against Monte-Carlo simulation of
Definition~\ref{def:ccsbm}, and every empirical quantity is computed directly from the
released data. For completeness we record the numerical choices.

Simulations draw $4\times10^{5}$ nodes per class per configuration, raised to
$1.5\times10^{6}$ for the crossover sweep, and errors are obtained by exhaustive
threshold search over the empirical distributions rather than by a plug-in formula, so
the reported values are attained minima. Edge-score laws are evaluated exactly from the
central and non-central $\chi^2$ distributions rather than by their Gaussian
approximation, which is used only where stated. The depth quantities require the
distance profile of an $L$-step walk on the $D$-regular tree; we obtain it by iterating
the birth--death recursion of Appendix~\ref{app:depth} to the required depth, which is
exact for the values of $L$ reported. Optimisations over a pruning operating point are
solved by bounded scalar minimisation on the logit of $\bar e_w$, and we verify that the
optimum is interior by re-solving on progressively wider domains; the reported values
are stable under that enlargement.

For the empirical networks, ties are treated as undirected and deduplicated, and
self-loops are removed. The camouflage ratio for an actor is the share of its labelled
alters belonging to the other class, and the reported $\hat\rho$ is the mean of that
share over strategic actors having at least one labelled tie; $\tilde D$ is the median
degree over the same set. Where a ratio of sums is compared against a prediction
derived from sums, both sides use the ratio-of-sums estimator, for the reason set out
in Section~\ref{sec:measuring}. The attribute separation $r$ is calibrated, for every dataset, from the
balanced error of a linear discriminant fitted on a random $70\%$ of accounts and
evaluated on the remainder, rather than from the moments of the attribute
distributions, for the reason given in Section~\ref{sec:r-estimate}. The decision
threshold is selected on the training split, not the evaluation split. Because the
estimate varies with the split, we average over repeated random splits and report the
spread.

All computations use double-precision floating point and standard scientific Python
libraries, and run to completion in a few minutes on a single CPU core.

\end{appendices}

\begin{thebibliography}{99}

\bibitem[Ajorlou et al.(2025)]{ajorlou2025dirichlet}
Ajorlou, H., Mateos, G., Ruiz, L., 2025. Dirichlet meets Horvitz and Thompson: estimating homophily in large networks via sampling. In: 2025 59th Asilomar Conference on Signals, Systems, and Computers. IEEE, pp. 1109-1113. doi:10.1109/IEEECONF67917.2025.11443313

\bibitem[An et al.(2026)]{an2026estimating}
An, W., Estrada, P., Estrada, J., Jacho-Chavez, D., 2026. Estimating peer influence in multilayer networks. Soc Networks.86:252-65. doi:10.1016/j.socnet.2026.03.003

\bibitem[Anderson et al.(1999)]{anderson1999interaction}
Anderson, B.S., Butts, C., Carley, K., 1999. The interaction of size and density with graph-level indices. Soc Networks.21(3):239-67. doi:10.1016/S0378-8733(99)00011-8

\bibitem[Ariyarathne et al.(2026)]{ariyarathne2026behavior}
Ariyarathne, I., Ariyarathne, G., Flammini, A., Menczer, F., Nwala, A., 2026. Behavior change as a signal for identifying social media manipulation. In: Proceedings of the 18th ACM Web Science Conference (WebSci '26). ACM, pp. 349-358. doi:10.1145/3795766.3799748

\bibitem[Baranwal et al.(2021)]{baranwal2021graph}
Baranwal, A., Fountoulakis, K., Jagannath, A., 2021. Graph convolution for semi-supervised classification: improved linear separability and out-of-distribution generalization. In: Proceedings of the 38th International Conference on Machine Learning (ICML 2021). PMLR.139:684-93.

\bibitem[Baranwal et al.(2023a)]{baranwal2023effects}
Baranwal, A., Fountoulakis, K., Jagannath, A., 2023a. Effects of graph convolutions in multi-layer networks. In: Proceedings of the 11th International Conference on Learning Representations (ICLR 2023).

\bibitem[Baranwal et al.(2023b)]{baranwal2023optimality}
Baranwal, A., Fountoulakis, K., Jagannath, A., 2023b. Optimality of message-passing architectures for sparse graphs. In: Advances in Neural Information Processing Systems 36 (NeurIPS 2023).

\bibitem[Baranwal(2026)]{baranwal2026depth}
Baranwal, A.R., 2026. The value of depth in message passing on sparse graphs: a Kesten--Stigum dichotomy. arXiv preprint arXiv:2607.16676.

\bibitem[Barone and Coscia(2018)]{barone2018birds}
Barone, M., Coscia, M., 2018. Birds of a feather scam together: trustworthiness homophily in a business network. Soc Networks.54:228-37. doi:10.1016/j.socnet.2018.01.009

\bibitem[Blanchard et al.(2017)]{blanchard2017machine}
Blanchard, P., El Mhamdi, E.M., Guerraoui, R., Stainer, J., 2017. Machine learning with adversaries: Byzantine tolerant gradient descent. In: Advances in Neural Information Processing Systems 30 (NeurIPS 2017). p. 118-28.

\bibitem[Bojanowski and Corten(2014)]{bojanowski2014measuring}
Bojanowski, M., Corten, R., 2014. Measuring segregation in social networks. Soc Networks.39:14-32. doi:10.1016/j.socnet.2014.04.001

\bibitem[Bojchevski and G\"unnemann(2019)]{bojchevski2019certifiable}
Bojchevski, A., G\"unnemann, S., 2019. Certifiable robustness to graph perturbations. In: Advances in Neural Information Processing Systems 32 (NeurIPS 2019).

\bibitem[Burt et al.(2021)]{burt2021network}
Burt, R.S., Reagans, R.E., Volvovsky, H.C., 2021. Network brokerage and the perception of leadership. Soc Networks.65:33-50. doi:10.1016/j.socnet.2020.09.002

\bibitem[Chandrasekhar et al.(2026)]{chandrasekhar2026multiplexing}
Chandrasekhar, A.G., Chaudhary, V., Golub, B., Jackson, M.O., 2026. Multiplexing in networks and diffusion. Proc. Natl. Acad. Sci. U.S.A. 123, e2534923123. doi:10.1073/pnas.2534923123

\bibitem[Chen et al.(2021)]{chen2021understanding}
Chen, L., Li, J., Peng, Q., Liu, Y., Zheng, Z., Yang, C., 2021. Understanding structural vulnerability in graph convolutional networks. In: Proceedings of the 30th International Joint Conference on Artificial Intelligence (IJCAI 2021). p. 2249-55. doi:10.24963/ijcai.2021/310

\bibitem[Choi and Young(2026)]{choi2026who}
Choi, E.C., Young, L.E., 2026. Who you know, what you believe, what you share: personal networks and health misinformation susceptibility. Soc Networks.87:163-72. doi:10.1016/j.socnet.2026.06.008

\bibitem[Cresci(2020)]{cresci2020decade}
Cresci, S., 2020. A decade of social bot detection. Commun ACM.63(10):72-83. doi:10.1145/3409116

\bibitem[Dang et al.(2026)]{dang2026xhbot}
Dang, Q.-V., Nguyen, P.-L., Le, D., Dinh, M.N., 2026. XHBot: eXplainable heterophily-aware graph neural networks for social bot detection. EAI Endorsed Trans AI Robotics [Internet]. 2026 Jul 7 [cited 2026 Jul 31];5. Available from: \url{https://publications.eai.eu/index.php/airo/article/view/12969}

\bibitem[Deshpande et al.(2018)]{deshpande2018contextual}
Deshpande, Y., Sen, S., Montanari, A., Mossel, E., 2018. Contextual stochastic block models. In: Advances in Neural Information Processing Systems 31 (NeurIPS 2018). p. 8590-602.

\bibitem[Dou et al.(2020)]{dou2020caregnn}
Dou, Y., Liu, Z., Sun, L., Deng, Y., Peng, H., Yu, P.S., 2020. Enhancing graph neural network-based fraud detectors against camouflaged fraudsters. In: Proceedings of the 29th ACM International Conference on Information and Knowledge Management (CIKM 2020). ACM. doi:10.1145/3340531.3411903

\bibitem[Duranthon and Zdeborov\'a(2025)]{duranthon2025statistical}
Duranthon, O., Zdeborov\'a, L., 2025. Statistical physics analysis of graph neural networks: approaching optimality in the contextual stochastic block model. Phys. Rev. X 15, 041026. doi:10.1103/lfxj-hbsk

\bibitem[Espinosa-Rada et al.(2026)]{espinosarada2026socioeconomic}
Espinosa-Rada, A., Bargsted, M., Ortiz Ruiz, F., 2026. Socioeconomic homogeneity in acquaintance networks: occupational prestige, education, and intergenerational mobility. Soc. Networks 86, 342-353. doi:10.1016/j.socnet.2026.04.012

\bibitem[Everett and Borgatti(2026)]{everett2026alter}
Everett, M.G., Borgatti, S.P., 2026. Alter composition with overlapping group memberships. Soc Networks.85:80-8. doi:10.1016/j.socnet.2025.12.001

\bibitem[Fluer et al.(2026)]{fluer2026survey}
Fluer, A., Laga, I., Graham, L., Almirol, E., Meyer, M., Schneider, J.A., Cummins, B., 2026. From survey data to social multiplex models: incorporating interlayer correlation from multiple data sources. Soc. Networks 86, 104-119. doi:10.1016/j.socnet.2026.01.005

\bibitem[Fountoulakis et al.(2023)]{fountoulakis2023graph}
Fountoulakis, K., Levi, A., Yang, S., Baranwal, A., Jagannath, A., 2023. Graph attention retrospective. J Mach Learn Res.24(246):1-52. Available from: \url{https://jmlr.org/papers/volume24/22-125/22-125.pdf}

\bibitem[Geisler et al.(2020)]{geisler2020reliable}
Geisler, S., Z\"ugner, D., G\"unnemann, S., 2020. Reliable graph neural networks via robust aggregation. In: Advances in Neural Information Processing Systems 33 (NeurIPS 2020).

\bibitem[Gosch et al.(2023)]{gosch2023revisiting}
Gosch, L., Geisler, S., Sturm, D., Charpentier, B., Z\"ugner, D., G\"unnemann, S., 2023. Revisiting robustness in graph machine learning. In: Proceedings of the 11th International Conference on Learning Representations (ICLR 2023).

\bibitem[He et al.(2025)]{he2025bothp}
He, B., Jiang, X., Wu, Q., Liu, H., Yang, Y., Liao, Y., 2025. Boosting bot detection via heterophily-aware representation learning and prototype-guided cluster discovery. In: Proceedings of the 31st ACM SIGKDD Conference on Knowledge Discovery and Data Mining (KDD 2025). ACM. doi:10.1145/3711896.3736862

\bibitem[Holland et al.(1983)]{holland1983stochastic}
Holland, P.W., Laskey, K.B., Leinhardt, S., 1983. Stochastic blockmodels: first steps. Soc Networks.5(2):109-37. doi:10.1016/0378-8733(83)90021-7

\bibitem[Jin et al.(2020)]{jin2020graph}
Jin, W., Ma, Y., Liu, X., Tang, X., Wang, S., Tang, J., 2020. Graph structure learning for robust graph neural networks. In: Proceedings of the 26th ACM SIGKDD International Conference on Knowledge Discovery and Data Mining (KDD 2020). ACM. p. 66-74. doi:10.1145/3394486.3403049

\bibitem[Kami\'nski et al.(2026)]{kaminski2026relationships}
Kami\'nski, B., Pra{\l}at, P., Wojnarowicz, A., Zawisza, M., 2026. Relationships between node degrees and hyperedge sizes in empirical hypergraphs. Soc. Networks 87, 119-134. doi:10.1016/j.socnet.2026.06.004

\bibitem[Krivitsky et al.(2009)]{krivitsky2009representing}
Krivitsky, P.N., Handcock, M.S., Raftery, A.E., Hoff, P.D., 2009. Representing degree distributions, clustering, and homophily in social networks with latent cluster random effects models. Soc Networks.31(3):204-13. doi:10.1016/j.socnet.2009.04.001

\bibitem[Lakdawala et al.(2026)]{lakdawala2026bonds}
Lakdawala, R., Leenders, R., Mulder, J., 2026. Not all bonds are created equal: dyadic latent class models for relational event data. Soc. Networks 87, 135-151. doi:10.1016/j.socnet.2026.06.006

\bibitem[Lazega and Pattison(1999)]{lazega1999multiplexity}
Lazega, E., Pattison, P.E., 1999. Multiplexity, generalized exchange and cooperation in organizations: a case study. Soc Networks.21(1):67-90. doi:10.1016/S0378-8733(99)00002-7

\bibitem[Leenders(2002)]{leenders2002modeling}
Leenders, R.T.h.A.J., 2002. Modeling social influence through network autocorrelation: constructing the weight matrix. Soc Networks.24(1):21-47. doi:10.1016/S0378-8733(01)00049-1

\bibitem[Li et al.(2018)]{li2018deeper}
Li, Q., Han, Z., Wu, X.M., 2018. Deeper insights into graph convolutional networks for semi-supervised learning. In: Proceedings of the 32nd AAAI Conference on Artificial Intelligence (AAAI 2018). p. 3538-45. doi:10.1609/aaai.v32i1.11604

\bibitem[Luan et al.(2022)]{luan2022revisiting}
Luan, S., Hua, C., Lu, Q., Zhu, J., Zhao, M., Zhang, S., et al., 2022. Revisiting heterophily for graph neural networks. In: Advances in Neural Information Processing Systems 35 (NeurIPS 2022).

\bibitem[Luan et al.(2023)]{luan2023when}
Luan, S., Hua, C., Lu, Q., Zhu, J., Zhao, M., Zhang, S., et al., 2023. When do graph neural networks help with node classification? Investigating the impact of homophily principle on node distinguishability. In: Advances in Neural Information Processing Systems 36 (NeurIPS 2023).

\bibitem[Ma et al.(2022)]{ma2022homophily}
Ma, Y., Liu, X., Shah, N., Tang, J., 2022. Is homophily a necessity for graph neural networks? In: International Conference on Learning Representations (ICLR 2022).

\bibitem[Ma et al.(2025)]{ma2025attention}
Ma, Z., Zhang, Q., Zhou, B., Zhang, Y., Hu, S., Wang, Z., 2025. Graph attention is not always beneficial: a theoretical analysis of graph attention mechanisms via contextual stochastic block models. In: Proceedings of the 42nd International Conference on Machine Learning (ICML 2025). arXiv:2412.15496.

\bibitem[Macdonald(2026)]{macdonald2026wiretap}
Macdonald, M., 2026. Wiretap surveillance of criminal networks does not resemble other types of collaborative networks. Soc. Networks 87, 12-23. doi:10.1016/j.socnet.2026.05.002

\bibitem[Mao et al.(2023)]{mao2023demystifying}
Mao, H., Chen, Z., Jin, W., Han, H., Ma, Y., Zhao, T., et al., 2023. Demystifying structural disparity in graph neural networks: can one size fit all? In: Advances in Neural Information Processing Systems 36 (NeurIPS 2023).

\bibitem[McMillan(2022)]{mcmillan2022worth}
McMillan, C., 2022. Worth the weight: conceptualizing and measuring strong versus weak tie homophily. Soc Networks.68:139-47. doi:10.1016/j.socnet.2021.06.003

\bibitem[Mujkanovic et al.(2022)]{mujkanovic2022are}
Mujkanovic, F., Geisler, S., G\"unnemann, S., Bojchevski, A., 2022. Are defenses for graph neural networks robust? In: Advances in Neural Information Processing Systems 35 (NeurIPS 2022).

\bibitem[Mukherjee et al.(2026)]{mukherjee2026optimal}
Mukherjee, K., Alom, Z., Ngo, T.G.B., Akcora, C.G., Kantarcioglu, M., 2026. Optimal transport-guided adversarial attacks on graph neural network-based bot detection. arXiv preprint arXiv:2602.00318. To appear, Proceedings of the 43rd International Conference on Machine Learning (ICML 2026).

\bibitem[Nguyen(2026)]{nguyen2026agency}
Nguyen, Q.T., 2026. Agency within structure: strategic behavior and network position in explaining premium buyer access in global buyer--seller networks. Soc. Networks 87, 188-209. doi:10.1016/j.socnet.2026.07.004

\bibitem[Noel and Nyhan(2011)]{noel2011unfriending}
Noel, H., Nyhan, B., 2011. The `unfriending' problem: the consequences of homophily in friendship retention for causal estimates of social influence. Soc Networks.33(3):211-8. doi:10.1016/j.socnet.2011.05.003

\bibitem[NT and Maehara(2019)]{nt2019revisiting}
NT, H., Maehara, T., 2019. Revisiting graph neural networks: all we have is low-pass filters. arXiv preprint. doi:10.48550/arXiv.1905.09550

\bibitem[Oono and Suzuki(2020)]{oono2020graph}
Oono, K., Suzuki, T., 2020. Graph neural networks exponentially lose expressive power for node classification. In: International Conference on Learning Representations (ICLR 2020).

\bibitem[Pan et al.(2026)]{pan2026camera}
Pan, J., Liu, Y., Zheng, Y., Chi, L., Liew, A.W.C., Pan, S., 2026. CAMERA: adapting to semantic camouflage in unsupervised text-attributed graph fraud detection. arXiv preprint arXiv:2605.20032. Accepted, IJCAI 2026.

\bibitem[Ran et al.(2025)]{ran2025identifying}
Ran, Y., Xiao, J., Xu, X.K., 2025. Identifying social bots via heterogeneous motifs based on na\"ive Bayes model. IEEE Trans. Inf. Forensics Secur. 20, 12754-12766. doi:10.1109/TIFS.2025.3635006

\bibitem[Roeling and Nicholls(2020)]{roeling2020imputation}
Roeling, M.P., Nicholls, G.K., 2020. Imputation of attributes in networked data using Bayesian autocorrelation regression models. Soc Networks.62:24-32. doi:10.1016/j.socnet.2020.02.005

\bibitem[Sch\"u{\ss}ler and Fuhse(2025)]{schussler2025duality}
Sch\"u{\ss}ler, A., Fuhse, J.A., 2025. The duality of network ties and attributes. Soc. Networks 83, 186-198. doi:10.1016/j.socnet.2025.07.002

\bibitem[Shea et al.(2019)]{shea2019cheater}
Shea, C.T., Lee, J., Menon, T., Im, D.-K., 2019. Cheater's hide and seek: strategic cognitive network activation during ethical decision making. Soc Networks.58:143-55. doi:10.1016/j.socnet.2019.03.005

\bibitem[Shi et al.(2025)]{shi2023mgtab}
Shi, S., Qiao, K., Liu, Z., Yang, J., Chen, C., Chen, J., et al., 2025. MGTAB: a multi-relational graph-based Twitter account detection benchmark. Neurocomputing.647:130490. doi:10.1016/j.neucom.2025.130490

\bibitem[Snijders et al.(2013)]{snijders2013model}
Snijders, T.A.B., Lomi, A., Torl\'o, V.J., 2013. A model for the multiplex dynamics of two-mode and one-mode networks, with an application to employment preference, friendship, and advice. Soc Networks.35(2):265-76. doi:10.1016/j.socnet.2012.05.005

\bibitem[Snijders(1981)]{snijders1981degree}
Snijders, T.A.B., 1981. The degree variance: an index of graph heterogeneity. Soc Networks.3(3):163-74. doi:10.1016/0378-8733(81)90014-9

\bibitem[Tang et al.(2026)]{tang2026twomode}
Tang, H., Mi, M., Ma, Y., Yang, H., Wang, F., 2026. Two-mode network autoregressive model for network analysis with core nodes. Soc. Networks 86, 75-87. doi:10.1016/j.socnet.2026.01.004

\bibitem[V\"or\"os and Snijders(2017)]{voros2017cluster}
V\"or\"os, A., Snijders, T.A.B., 2017. Cluster analysis of multiplex networks: defining composite network measures. Soc Networks.49:93-112. doi:10.1016/j.socnet.2017.01.002

\bibitem[V\"or\"os et al.(2019)]{voros2019limits}
V\"or\"os, A., Block, P., Boda, Z., 2019. Limits to inferring status from friendship relations. Soc Networks.59:77-97. doi:10.1016/j.socnet.2019.05.007

\bibitem[Weiss(2026)]{weiss2026cost}
Weiss, E., 2026. Cost-sensitive neighborhood aggregation for heterophilous graphs: when does per-edge routing help? arXiv preprint. doi:10.48550/arXiv.2603.24291

\bibitem[Yin et al.(2018)]{yin2018byzantine}
Yin, D., Chen, Y., Kannan, R., Bartlett, P., 2018. Byzantine-robust distributed learning: towards optimal statistical rates. In: Proceedings of the 35th International Conference on Machine Learning (ICML 2018). PMLR 80. p. 5650-9.

\bibitem[Yu et al.(2018)]{yu2018detecting}
Yu, L., Woodall, W.H., Tsui, K.-L., 2018. Detecting node propensity changes in the dynamic degree corrected stochastic block model. Soc Networks.54:209-27. doi:10.1016/j.socnet.2018.03.004

\bibitem[Z\"ugner and G\"unnemann(2019)]{zugner2019adversarial}
Z\"ugner, D., G\"unnemann, S., 2019. Adversarial attacks on graph neural networks via meta learning. In: Proceedings of the 7th International Conference on Learning Representations (ICLR 2019).

\bibitem[Z\"ugner et al.(2018)]{zugner2018adversarial}
Z\"ugner, D., Akbarnejad, A., G\"unnemann, S., 2018. Adversarial attacks on neural networks for graph data. In: Proceedings of the 24th ACM SIGKDD International Conference on Knowledge Discovery and Data Mining (KDD 2018). ACM. p. 2847-56. doi:10.1145/3219819.3220078

\bibitem[Zhang and Zitnik(2020)]{zhang2020gnnguard}
Zhang, X., Zitnik, M., 2020. GNNGuard: defending graph neural networks against adversarial attacks. In: Advances in Neural Information Processing Systems 33 (NeurIPS 2020).

\bibitem[Zhu et al.(2020)]{zhu2020beyond}
Zhu, J., Yan, Y., Zhao, L., Heimann, M., Akoglu, L., Koutra, D., 2020. Beyond homophily in graph neural networks: current limitations and effective designs. In: Advances in Neural Information Processing Systems 33 (NeurIPS 2020).

\bibitem[Zhu et al.(2022)]{zhu2022how}
Zhu, J., Jin, J., Loveland, D., Schaub, M.T., Koutra, D., 2022. How does heterophily impact the robustness of graph neural networks? Theoretical connections and practical implications. In: Proceedings of the 28th ACM SIGKDD Conference on Knowledge Discovery and Data Mining (KDD 2022). ACM. p. 2637-47. doi:10.1145/3534678.3539418

\end{thebibliography}

\end{document}
